% Generated mechanically from the frozen spaces-groupoids.tex target.
% Do not edit this generated file; edit the bound locale source instead.

% OK, start here.
%

\setcounter{chapter}{77}
\chapter{代数空间中的群胚}



\phantomsection
\label{spaces-groupoids-section-phantom}


\section{引言}
\label{spaces-groupoids-section-introduction}

\noindent
本章专论代数空间中群胚的一般理论。我们推荐阅读 Keel 和 Mori 的优美论文
\cite{K-M}。

\medskip\noindent
这里的很多内容重复了群胚概形一章中已经讨论过的内容，见《群胚》第
\ref{groupoids-section-introduction} 节。这里关于商栈的讨论则是新的。


\section{约定}
\label{spaces-groupoids-section-conventions}

\noindent
我们始终假设所有概形都包含在一个大 fppf 位形 $\Sch_{fppf}$ 中，
并且所有考虑的环 $A$ 都具有如下性质：$\Spec(A)$（在同构意义下）
是这个大位形的对象。

\medskip\noindent
设 $S$ 为概形，$X$ 为 $S$ 上的代数空间。在本章及下一章中，我们用
$X \times_S X$ 表示 $X$ 与自身的积（取在 $S$ 上代数空间的范畴中），
而不用 $X \times X$。

\medskip\noindent
我们继续沿用从指标 $0$ 开始标记投影映射的约定，因而有
$\text{pr}_0 : X \times_S Y \to X$ 和
$\text{pr}_1 : X \times_S Y \to Y$。




\section{记号}
\label{spaces-groupoids-section-notation}

\noindent
设 $S$ 为概形；它将是我们的基概形，所有代数空间都取在 $S$ 上。
设 $B$ 为 $S$ 上的代数空间；它将是我们的基代数空间，其他代数空间以及
概形往往取在 $B$ 上。若称 $X$ 是 $B$ 上的代数空间，意思是 $X$ 是
$S$ 上的代数空间，并配备结构态射 $X \to B$。此外，我们尽量保留字母
$T$ 来表示 $B$ 上的“测试”概形。换言之，$T$ 是配备结构态射
$T \to B$ 的概形。在这种情形下，我们用 $X(T)$ 表示 $T$-值点所成的集合；
这些点取值于 $X$ 且位于 $B$ 上。用公式写就是
$$
X(T) = \Mor_B(T, X).
$$
类似地，给定另一个代数空间 $Y$（在 $B$ 上），置
$$
X(Y) = \Mor_B(Y, X).
$$
假设如上给定代数空间 $X$、$Y$（在 $B$ 上），以及态射
$f : X \to Y$（在 $B$ 上）。对任意概形 $T$（在 $B$ 上），得到诱导的集合映射
$$
f : X(T) \longrightarrow Y(T)
$$
而且它对概形 $T$（在 $B$ 上）具有函子性。由于 $f$ 是
$(\Sch/S)_{fppf}$ 上、位于层 $B$ 之上的层映射，显然 $f$ 决定这个规则，
也由这个规则决定。更一般地，我们对纤维积之间的映射采用相同记号。
例如，若 $X$、$Y$、$Z$ 是 $B$ 上的代数空间，而
$m : X \times_B Y \to Z \times_B Z$ 是 $B$ 上代数空间的态射，
那么可把 $m$ 看作对应于如下 $T$-值点之间的一族映射：
$$
X(T) \times Y(T) \longrightarrow Z(T) \times Z(T).
$$
其余情形依此类推。

\medskip\noindent
最后，给定代数空间的两个映射 $f, g : X \to Y$（在 $B$ 上），若诱导映射
$f, g : X(T) \to Y(T)$ 对每个概形 $T$（在 $B$ 上）都相等，那么
$f = g$；因而映射 $f, g : X(Z) \to Y(Z)$ 对每个第三个代数空间
$Z$（在 $B$ 上）也相等。例如，要检验群代数空间 $G$（在 $B$ 上）的
% P09-ERR-0866
公理，只须在 $T$-值点上检验图表的交换性；这里 $T$ 是 $B$ 上的概形，
% P09-ERR-0867
具体做法如下面的定义 \ref{spaces-groupoids-definition-group-space} 所示。


\section{等价关系}
\label{spaces-groupoids-section-equivalence-relations}

\noindent
记号参见《群胚》第 \ref{groupoids-section-equivalence-relations} 节。

\begin{definition}
\label{spaces-groupoids-definition-equivalence-relation}
设 $B \to S$ 如第 \ref{spaces-groupoids-section-notation} 节所示。设 $U$ 为代数空间（在
$B$ 上）。
\begin{enumerate}
\item $U$ 在 $B$ 上的一个\emph{预关系}，是任意态射
$j : R \to U \times_B U$（它是 $B$ 上代数空间的态射）。在这种情形下，置
$t = \text{pr}_0 \circ j$ 和 $s = \text{pr}_1 \circ j$，从而
$j = (t, s)$。
\item $U$ 在 $B$ 上的一个\emph{关系}，是单态射
$j : R \to U \times_B U$（它是 $B$ 上代数空间的态射）。
\item 一个\emph{预等价关系}，是预关系 $j : R \to U \times_B U$，
它满足：映射 $j : R(T) \to U(T) \times U(T)$ 的像对所有概形
$T$（在 $B$ 上）都是等价关系。
\item 称态射 $R \to U \times_B U$（它是 $B$ 上代数空间的态射）是
\emph{$U$ 在 $B$ 上的等价关系}，当且仅当对每个 $T$（在 $B$ 上），
相应的 $T$-值点（属于 $R$）在相应的 $T$-值点集（属于 $U$）上定义一个
等价关系。
\end{enumerate}
\end{definition}

\noindent
换言之，等价关系就是满足 $j$ 为关系的预等价关系。

\begin{lemma}
\label{spaces-groupoids-lemma-restrict-relation}
设 $B \to S$ 如第 \ref{spaces-groupoids-section-notation} 节所示。设 $U$ 为代数空间（在
$B$ 上）。设 $j : R \to U \times_B U$ 为预关系。设
$g : U' \to U$ 为代数空间的态射（在 $B$ 上）。最后，置
$$
R' = (U' \times_B U')\times_{U \times_B U} R
\xrightarrow{j'}
U' \times_B U'
$$
那么 $j'$ 是 $U'$ 在 $B$ 上的预关系。如果 $j$ 是关系，那么 $j'$ 也是关系。
如果 $j$ 是预等价关系，那么 $j'$ 也是预等价关系。如果 $j$ 是等价关系，
那么 $j'$ 也是等价关系。
\end{lemma}

\begin{proof}
略。
\end{proof}







\begin{definition}
\label{spaces-groupoids-definition-restrict-relation}
设 $B \to S$ 如第 \ref{spaces-groupoids-section-notation} 节所示。设 $U$ 为代数空间（在
$B$ 上）。设 $j : R \to U \times_B U$ 为预关系。设
$g : U' \to U$ 为代数空间的态射（在 $B$ 上）。引理
\ref{spaces-groupoids-lemma-restrict-relation} 中的预关系
$j' : R' \to U' \times_B U'$ 称为预关系 $j$ 到 $U'$ 的\emph{限制}或
\emph{拉回}。在这种情形下，有时写成 $R' = R|_{U'}$。
\end{definition}

\begin{lemma}
\label{spaces-groupoids-lemma-pre-equivalence-equivalence-relation-points}
设 $B \to S$ 如第 \ref{spaces-groupoids-section-notation} 节所示。设
$j : R \to U \times_B U$ 为代数空间的预关系（在 $B$ 上）。考虑由下述规则
在 $|U|$ 上定义的关系：
$$
x \sim y
\Leftrightarrow
\exists\ r \in |R| :
t(r) = x,
s(r) = y.
$$
如果 $j$ 是预等价关系，那么这就是等价关系。
\end{lemma}

\begin{proof}
假设 $x \sim y$ 且 $y \sim z$。选取 $r \in |R|$，使得
$t(r) = x$、$s(r) = y$；再选取 $r' \in |R|$，使得
$t(r') = y$、$s(r') = z$。可以选取域 $K$，使得 $r$ 和 $r'$ 可由态射
$r, r' : \Spec(K) \to R$ 表示，并且 $s \circ r = t \circ r'$。
记 $x = t \circ r$、$y = s \circ r = t \circ r'$ 和
$z = s \circ r'$，于是 $x, y, z : \Spec(K) \to U$。由构造，
$(x, y) \in j(R(K))$ 且 $(y, z) \in j(R(K))$。由于 $j$ 是预等价关系，
可知 $(x, z) \in j(R(K))$ 也成立。这显然推出 $x \sim z$。

\medskip\noindent
略去 $\sim$ 具有自反性和对称性的证明。
\end{proof}
\section{群代数空间}
\label{spaces-groupoids-section-group-spaces}

\noindent
记号参见《群胚》第 \ref{groupoids-section-group-schemes} 节。

\begin{definition}
\label{spaces-groupoids-definition-group-space}
设 $B \to S$ 如第 \ref{spaces-groupoids-section-notation} 节所示。
\begin{enumerate}
\item 一个\emph{$B$ 上的群代数空间}是偶对 $(G, m)$，其中 $G$ 是
$B$ 上的代数空间，而 $m : G \times_B G \to G$ 是 $B$ 上代数空间的态射，
并具有如下性质：对每个概形 $T$（在 $B$ 上），偶对 $(G(T), m)$ 是群。
\item 一个\emph{态射 $\psi : (G, m) \to (G', m')$（它是 $B$ 上
群代数空间的态射）}，是态射 $\psi : G \to G'$（它是 $B$ 上代数空间的态射），
并且对每个 $T/B$，诱导映射
$\psi : G(T) \to G'(T)$ 都是群同态。
\end{enumerate}
\end{definition}

\noindent
设 $(G, m)$ 为代数空间 $B$ 上的群代数空间。根据《群胚》第
\ref{groupoids-section-group-schemes} 节的讨论，得到 $B$ 上代数空间的态射
（单位元）$e : B \to G$ 和（逆元）$i : G \to G$，使得对每个 $T$，
四元组 $(G(T), m, e, i)$ 满足群公理。

\medskip\noindent
设 $(G, m)$、$(G', m')$ 为 $B$ 上的群代数空间。设
$f : G \to G'$ 为 $B$ 上代数空间的态射。由定义，$f$ 是 $B$ 上群代数空间的
态射，当且仅当下图交换：
$$
\xymatrix{
G \times_B G \ar[r]_-{f \times f} \ar[d]_m &
% P09-ERR-0868
G' \times_B G' \ar[d]^m \\
G \ar[r]^f & G'
}
$$

\begin{lemma}
\label{spaces-groupoids-lemma-base-change-group-space}
设 $B \to S$ 如第 \ref{spaces-groupoids-section-notation} 节所示。设 $(G, m)$ 为 $B$ 上的
群代数空间。设 $B' \to B$ 为代数空间的态射。拉回
$(G_{B'}, m_{B'})$ 是 $B'$ 上的群代数空间。
\end{lemma}

\begin{proof}
略。
\end{proof}







\section{群代数空间的性质}
\label{spaces-groupoids-section-properties-group-spaces}

\noindent
本节汇集群代数空间在任意基底上成立的一些简单性质。

\begin{lemma}
\label{spaces-groupoids-lemma-group-scheme-separated}
设 $S$ 为概形，$B$ 为 $S$ 上的代数空间，$G$ 为 $B$ 上的群代数空间。
那么 $G \to B$ 分离（相应地，拟分离；局部分离），当且仅当单位元态射
$e : B \to G$ 是闭浸入（相应地，拟紧；浸入）。
\end{lemma}

\begin{proof}
回忆《空间的态射》引理 \ref{spaces-morphisms-lemma-section-immersion}：
$e$ 在 $G \to B$ 分离（相应地，拟分离；局部分离）时是闭浸入
（相应地，拟紧；浸入）。为证明逆命题，考虑图表
$$
\xymatrix{
G \ar[r]_-{\Delta_{G/B}} \ar[d] &
G \times_B G \ar[d]^{(g, g') \mapsto m(i(g), g')} \\
B \ar[r]^e & G
}
$$
运用代数几何中的函子点观点，不难验证这个图表是笛卡尔图表。换言之，
$\Delta_{G/B}$ 是 $e$ 的基变换。因此，如果 $e$ 是闭浸入（相应地，拟紧；
浸入），那么 $\Delta_{G/B}$ 也具有相应性质；参见《空间》引理
\ref{spaces-lemma-base-change-immersions}（相应地，《空间的态射》引理
\ref{spaces-morphisms-lemma-base-change-quasi-compact}；《空间》引理
\ref{spaces-lemma-base-change-immersions}）。
\end{proof}

\begin{lemma}
\label{spaces-groupoids-lemma-group-scheme-unramified-or-lqf}
设 $S$ 为概形，$B$ 为 $S$ 上的代数空间，$G$ 为 $B$ 上的群代数空间。
假设 $G \to B$ 局部有限型。那么 $G \to B$ 非分歧（相应地，局部拟有限），
当且仅当 $G \to B$ 在 $e(b)$ 处非分歧（相应地，拟有限）；这对所有
$b \in |B|$ 都成立。
\end{lemma}

\begin{proof}
由《空间的态射》引理 \ref{spaces-morphisms-lemma-where-unramified}
（相应地，《空间的态射》引理
\ref{spaces-morphisms-lemma-base-change-quasi-finite-locus}），存在最大的
开子空间 $U \subset G$，使得 $U \to B$ 非分歧（相应地，局部拟有限），
且 $U$ 的构造与基变换交换。因此，可归约到 $B = \Spec(k)$ 是域的谱这一情形。
设 $g \in G(K)$ 为取值于扩域 $K/k$ 的点。为了检验 $g$ 是否属于 $U$，
可以把基底变换到 $K$。所以只须证明
$$
G \to \Spec(k)\text{ 非分歧于 }e
\Leftrightarrow
G \to \Spec(k)\text{ 非分歧于 }g
$$
其中在 $k$ 上，$g$ 是有理点（相应地，在 $g$ 与 $e$ 处拟有限的情形也一样）。
由于沿 $g$ 的平移是
$G$ 在 $k$ 上的自同构，这一点很清楚。
\end{proof}

\begin{lemma}
\label{spaces-groupoids-lemma-open-over-which-unramified-or-lqf}
设 $S$ 为概形，$B$ 为 $S$ 上的代数空间，$G$ 为 $B$ 上的群代数空间。
假设 $G \to B$ 局部有限型。
\begin{enumerate}
\item 存在最大的开子空间 $U \subset B$，使得 $G_U \to U$ 非分歧，
且 $U$ 的构造与基变换交换。
\item 存在最大的开子空间 $U \subset B$，使得 $G_U \to U$ 局部拟有限，
且 $U$ 的构造与基变换交换。
\end{enumerate}
\end{lemma}

\begin{proof}
由《空间的态射》引理 \ref{spaces-morphisms-lemma-where-unramified}
（相应地，《空间的态射》引理
\ref{spaces-morphisms-lemma-base-change-quasi-finite-locus}），存在最大的
开子空间 $W \subset G$，使得 $W \to B$ 非分歧（相应地，局部拟有限）。
此外，$W$ 的构造与基变换交换。由引理
\ref{spaces-groupoids-lemma-group-scheme-unramified-or-lqf}，可知两种情形下都有
$U = e^{-1}(W)$。
\end{proof}
\section{群代数空间的例}
\label{spaces-groupoids-section-examples-group-spaces}

\noindent
如果 $G \to S$ 是基概形 $S$ 上的群概形，那么由引理
\ref{spaces-groupoids-lemma-base-change-group-space}，基变换 $G_B$ 是群代数空间；这里 $B$ 是
$S$ 上任意代数空间，而所得群代数空间取在 $B$ 上。
% P09-ERR-0869
我们将在下面的例中频繁使用这一点。

\begin{example}[乘法群代数空间]
\label{spaces-groupoids-example-multiplicative-group}
设 $B \to S$ 如第 \ref{spaces-groupoids-section-notation} 节所示。考虑如下函子：它把任意概形
$T$（在 $B$ 上）映到结构层整体截面中单位元所成的群
$\Gamma(T, \mathcal{O}_T^*)$。这个函子由群代数空间
$$
\mathbf{G}_{m, B} = B \times_S \mathbf{G}_{m, S}
$$
表示；该空间取在 $B$ 上。这里 $\mathbf{G}_{m, S}$ 是 $S$ 上的乘法群概形，见《群胚》例
\ref{groupoids-example-multiplicative-group}。
\end{example}

\begin{example}[作为群代数空间的单位根]
\label{spaces-groupoids-example-roots-of-unity}
设 $B \to S$ 如第 \ref{spaces-groupoids-section-notation} 节所示。设 $n \in \mathbf{N}$。
考虑如下函子：它把任意概形 $T$（在 $B$ 上）映到
$\Gamma(T, \mathcal{O}_T^*)$ 中由 $n$ 次单位根组成的子群。这个函子由
群代数空间
$$
\mu_{n, B} = B \times_S \mu_{n, S}
$$
表示；该空间取在 $B$ 上。这里 $\mu_{n, S}$ 是 $n$ 次单位根群概形（在 $S$ 上），见《群胚》例
\ref{groupoids-example-roots-of-unity}。
\end{example}

\begin{example}[加法群代数空间]
\label{spaces-groupoids-example-additive-group}
设 $B \to S$ 如第 \ref{spaces-groupoids-section-notation} 节所示。考虑如下函子：它把任意概形
$T$（在 $B$ 上）映到结构层整体截面所成的群
$\Gamma(T, \mathcal{O}_T)$。这个函子由群代数空间
$$
\mathbf{G}_{a, B} = B \times_S \mathbf{G}_{a, S}
$$
表示；该空间取在 $B$ 上。这里 $\mathbf{G}_{a, S}$ 是 $S$ 上的加法群概形，见《群胚》例
\ref{groupoids-example-additive-group}。
\end{example}

\begin{example}[一般线性群代数空间]
\label{spaces-groupoids-example-general-linear-group}
设 $B \to S$ 如第 \ref{spaces-groupoids-section-notation} 节所示。设 $n \geq 1$。
考虑如下函子：它把任意概形 $T$（在 $B$ 上）映到群
$$
\text{GL}_n(\Gamma(T, \mathcal{O}_T))
$$
即结构层整体截面上的可逆 $n \times n$ 矩阵群。这个函子由群代数空间
$$
\text{GL}_{n, B} = B \times_S \text{GL}_{n, S}
$$
表示；该空间取在 $B$ 上。这里
% P09-ERR-0870
$\mathbf{G}_{m, S}$ 是 $S$ 上的一般线性群概形，见《群胚》例
\ref{groupoids-example-general-linear-group}。
\end{example}

\begin{example}
\label{spaces-groupoids-example-determinant}
设 $B \to S$ 如第 \ref{spaces-groupoids-section-notation} 节所示。设 $n \geq 1$。
行列式定义群代数空间的态射
$$
\det : \text{GL}_{n, B} \longrightarrow \mathbf{G}_{m, B}
$$
这是 $B$ 上的态射。它是《群胚》例
\ref{groupoids-example-determinant} 中 $S$ 上行列式态射的基变换。
\end{example}

\begin{example}[常值群代数空间]
\label{spaces-groupoids-example-constant-group}
设 $B \to S$ 如第 \ref{spaces-groupoids-section-notation} 节所示。设 $G$ 为抽象群。
考虑如下函子：它把任意概形 $T$（在 $B$ 上）映到局部常值映射
$T \to G$ 所成的群（这里 $T$ 取 Zariski 拓扑，$G$ 取离散拓扑）。
这个函子由群代数空间
$$
G_B = B \times_S G_S
$$
表示；该空间取在 $B$ 上。这里 $G_S$ 是《群胚》例 \ref{groupoids-example-constant-group}
中引入的常值群概形。
\end{example}
\section{群代数空间的作用}
\label{spaces-groupoids-section-action-group-space}

\noindent
记号参见《群胚》第 \ref{groupoids-section-action-group-scheme} 节。

\begin{definition}
\label{spaces-groupoids-definition-action-group-space}
设 $B \to S$ 如第 \ref{spaces-groupoids-section-notation} 节所示。设 $(G, m)$ 为 $B$ 上的
群代数空间。设 $X$ 为 $B$ 上的代数空间。
\begin{enumerate}
\item \emph{$G$ 在代数空间 $X/B$ 上的作用}，是态射
$a : G \times_B X \to X$（在 $B$ 上），它满足：对每个概形 $T$（在 $B$ 上），映射
$a : G(T) \times X(T) \to X(T)$ 定义一个 $G(T)$-集结构，其底层集合为
$X(T)$。
\item 假设 $X$、$Y$ 是 $B$ 上的代数空间，并且各自配备 $G$ 的作用。
一个\emph{等变}态射，更精确地说，一个\emph{$G$-等变态射}
$\psi : X \to Y$，是 $B$ 上代数空间的态射，它满足：对每个 $T$（在
$B$ 上），映射 $\psi : X(T) \to Y(T)$ 是 $G(T)$-集的态射。
\end{enumerate}
\end{definition}

\noindent
在情形 (1) 中，这意味着图表
\begin{equation}
\label{spaces-groupoids-equation-action}
\xymatrix{
G \times_B G \times_B X \ar[r]_-{1_G \times a} \ar[d]_{m \times 1_X} &
G \times_B X \ar[d]^a \\
G \times_B X \ar[r]^a & X
}
\quad
\xymatrix{
G \times_B X \ar[r]_-a & X \\
X\ar[u]^{e \times 1_X} \ar[ru]_{1_X}
}
\end{equation}
交换。在情形 (2) 中，这仅意味着图表
$$
\xymatrix{
% P09-ERR-0871
G \times_B X \ar[r]_-{\text{id} \times f} \ar[d]_a &
G \times_B Y \ar[d]^a \\
X \ar[r]^f & Y
}
$$
交换。

\begin{definition}
\label{spaces-groupoids-definition-free-action}
设 $B \to S$、$G \to B$ 和 $X \to B$ 如定义
\ref{spaces-groupoids-definition-action-group-space} 所示。设
$a : G \times_B X \to X$ 为 $G$ 在 $X/B$ 上的作用。如果对每个概形
$T$（在 $B$ 上），作用 $a : G(T) \times X(T) \to X(T)$ 都是群
$G(T)$ 在集合 $X(T)$ 上的自由作用，就称该作用\emph{自由}。
\end{definition}

\begin{lemma}
\label{spaces-groupoids-lemma-free-action}
情形如定义 \ref{spaces-groupoids-definition-free-action} 所示。
% P09-ERR-0872
作用 $a$ 自由，当且仅当
$$
G \times_B X \to X \times_B X, \quad (g, x) \mapsto (a(g, x), x)
$$
是代数空间的单态射。
\end{lemma}

\begin{proof}
由定义立即得到。
\end{proof}
\section{主齐性空间}
\label{spaces-groupoids-section-principal-homogeneous}

\noindent
本节对应于《群胚》第
\ref{groupoids-section-principal-homogeneous} 节。我们建议先阅读该节。

\begin{definition}
\label{spaces-groupoids-definition-pseudo-torsor}
设 $S$ 为概形。设 $B$ 为 $S$ 上的代数空间。设 $(G, m)$ 为 $B$ 上的
群代数空间。设 $X$ 为 $B$ 上的代数空间，并设
$a : G \times_B X \to X$ 为 $G$ 在 $X$ 上的作用。
\begin{enumerate}
\item 称 $X$ 为一个\emph{伪 $G$-挠子}，或称 $X$
\emph{在 $G$ 作用下形式主齐性}，如果诱导态射
$G \times_B X \to X \times_B X$，
$(g, x) \mapsto (a(g, x), x)$ 是同构。
\item 一个伪 $G$-挠子 $X$ 称为\emph{平凡的}，如果存在
% P09-ERR-0873
一个 $G$-等变同构 $G \to X$（在 $B$ 上），其中 $G$ 通过左乘作用在
$G$ 上。
\end{enumerate}
\end{definition}

\noindent
显然，如果 $B' \to B$ 是代数空间的态射，那么伪挠子的拉回
$X_{B'}$（原挠子为伪 $G$-挠子，取在 $B$ 上）是伪
$G_{B'}$-挠子（取在 $B'$ 上）。

\begin{lemma}
\label{spaces-groupoids-lemma-characterize-trivial-pseudo-torsors}
情形如定义 \ref{spaces-groupoids-definition-pseudo-torsor} 所示。
\begin{enumerate}
\item 代数空间 $X$ 是伪 $G$-挠子，当且仅当对每个概形 $T$（在 $B$
上），集合 $X(T)$ 或者为空，或者群 $G(T)$ 在 $X(T)$ 上的作用单传递。
\item 伪 $G$-挠子 $X$ 平凡，当且仅当态射 $X \to B$ 有截面。
\end{enumerate}
\end{lemma}

\begin{proof}
略。
\end{proof}

\begin{definition}
\label{spaces-groupoids-definition-principal-homogeneous-space}
设 $S$ 为概形。设 $B$ 为 $S$ 上的代数空间。设 $(G, m)$ 为 $B$ 上的
群代数空间。设 $X$ 为伪 $G$-挠子（在 $B$ 上）。
\begin{enumerate}
\item 称 $X$ 为一个\emph{主齐性空间}，更精确地说，称它为一个
\emph{主齐性 $G$-空间（在 $B$ 上）}，如果存在一个 fpqc 覆盖\footnote{《群胚》定义
\ref{groupoids-definition-principal-homogeneous-space} 中默认的挠子类型，
是在某个 fpqc 覆盖上平凡的伪挠子。由于
% P09-ERR-0874
$G$ 作为代数空间可以视为群层，那里已经有一个 $G$-挠子的概念；这个概念
% P09-ERR-0875
对应于 fppf-挠子，见引理 \ref{spaces-groupoids-lemma-torsor}。因此，对于 fpqc 局部平凡
的伪挠子，我们使用“主齐性空间”这一名称，并尽量避免在这种情形下使用
“挠子”一词。}
$\{B_i \to B\}_{i \in I}$，使得每个 $X_{B_i} \to B_i$ 都有截面
（即它是平凡的伪 $G_{B_i}$-挠子）。
\item 设 $\tau \in \{Zariski, \etale, smooth, syntomic, fppf\}$。
称 $X$ 为一个\emph{$G$-挠子（取 $\tau$ 拓扑）}，或称为一个
\emph{$\tau$ $G$-挠子}，或简称为一个\emph{$\tau$-挠子}，如果存在
一个 $\tau$ 覆盖 $\{B_i \to B\}_{i \in I}$，使得每个
$X_{B_i} \to B_i$ 都有截面。
\item 如果 $X$ 是主齐性 $G$-空间（在 $B$ 上），并且它是
\'etale 拓扑中的挠子，就称它\emph{拟等平凡}。
\item 如果 $X$ 是主齐性 $G$-空间（在 $B$ 上），并且它是 Zariski
拓扑中的挠子，就称它\emph{局部平凡}。
\end{enumerate}
\end{definition}

\noindent
我们有时说“设 $X$ 为 $G$-主齐性空间（在 $B$ 上）”，意思是说 $X$ 是
$B$ 上配备了 $G$ 的一个作用的代数空间，并且这个作用使它成为 $B$ 上的
主齐性空间。下面我们证明，在两种记号都适用时，它们相互一致。

\begin{lemma}
\label{spaces-groupoids-lemma-torsor}
设 $S$ 为概形。设 $(G, m)$ 为 $S$ 上的群代数空间。设 $X$ 为 $S$
上的代数空间，并设 $a : G \times_S X \to X$ 为 $G$ 在 $X$ 上的
作用。那么，按定义 \ref{spaces-groupoids-definition-principal-homogeneous-space}，
$X$ 是 $G$-挠子（取 $fppf$-拓扑），当且仅当按《位点上的上同调》定义
\ref{sites-cohomology-definition-torsor}，$X$ 是 $G$-挠子（在
$(\Sch/S)_{fppf}$ 上）。
\end{lemma}

\begin{proof}
略。
\end{proof}

\begin{lemma}
\label{spaces-groupoids-lemma-pseudo-torsor-implications}
设 $S$ 为概形。设 $B$ 为 $S$ 上的代数空间。设 $G$ 为 $B$ 上的群
代数空间。设 $X$ 为伪 $G$-挠子（在 $B$ 上）。假设
% P09-ERR-0876
$G$ 和 $X$ 在 $B$ 上局部有限型。
\begin{enumerate}
\item 如果 $G \to B$ 非分歧，那么 $X \to B$ 非分歧。
\item 如果 $G \to B$ 局部拟有限，那么 $X \to B$ 局部拟有限。
\end{enumerate}
\end{lemma}

\begin{proof}
(1) 的证明。由《空间的态射》引理
\ref{spaces-morphisms-lemma-where-unramified}，我们约化到 $B$ 是域的
谱的情形。如果 $X$ 为空，则结论成立。如果 $X$ 非空，那么在扩大该域之后，
可以假设 $X$ 有一个点。于是 $G \cong X$，结论成立。

\medskip\noindent
(2) 的证明完全相同，只需使用《空间的态射》引理
\ref{spaces-morphisms-lemma-base-change-quasi-finite-locus}。
\end{proof}
\section{等变拟凝聚层}
\label{spaces-groupoids-section-equivariant}

\noindent
请与《群胚》第 \ref{groupoids-section-equivariant} 节比较。

\begin{definition}
\label{spaces-groupoids-definition-equivariant-module}
设 $B \to S$ 如第 \ref{spaces-groupoids-section-notation} 节所示。设 $(G, m)$ 为 $B$ 上的
群代数空间，并设 $a : G \times_B X \to X$ 为 $G$ 在代数空间 $X$ 上的
作用，该空间取在 $B$ 上。一个\emph{$G$-等变拟凝聚
$\mathcal{O}_X$-模}，或简称一个\emph{等变拟凝聚 $\mathcal{O}_X$-模}，
是一个偶 $(\mathcal{F}, \alpha)$，其中 $\mathcal{F}$ 是拟凝聚
$\mathcal{O}_X$-模，而 $\alpha$ 是 $\mathcal{O}_{G \times_B X}$-模映射
$$
\alpha : a^*\mathcal{F} \longrightarrow \text{pr}_1^*\mathcal{F}
$$
这里 $\text{pr}_1 : G \times_B X \to X$ 是投影，并且满足：
\begin{enumerate}
\item 图表
$$
\xymatrix{
% P09-ERR-0877
(1_G \times a)^*\text{pr}_2^*\mathcal{F} \ar[r]_-{\text{pr}_{12}^*\alpha} &
\text{pr}_2^*\mathcal{F} \\
(1_G \times a)^*a^*\mathcal{F} \ar[u]^{(1_G \times a)^*\alpha} \ar@{=}[r] &
(m \times 1_X)^*a^*\mathcal{F} \ar[u]_{(m \times 1_X)^*\alpha}
}
$$
% P09-ERR-0878
在 $\mathcal{O}_{G \times_B G \times_B X}$-模范畴中交换；
\item 拉回
$$
(e \times 1_X)^*\alpha : \mathcal{F} \longrightarrow \mathcal{F}
$$
是恒等映射。
\end{enumerate}
解释可与式 (\ref{spaces-groupoids-equation-action}) 的有关图表比较。
\end{definition}

\noindent
注意，第一个图表的交换性保证 $(e \times 1_X)^*\alpha$ 是
$\mathcal{F}$ 上的幂等算子；因此条件 (2) 恰好是说它为同构。

\begin{lemma}
\label{spaces-groupoids-lemma-pullback-equivariant}
设 $B \to S$ 如第 \ref{spaces-groupoids-section-notation} 节所示。设 $G$ 为 $B$ 上的
群代数空间。设 $f : X \to Y$ 为一个 $G$-等变态射，它连接两个在 $B$
上且配备 $G$-作用的代数空间。
% P09-ERR-0879
拉回 $f^*$ 由
$(\mathcal{F}, \alpha) \mapsto (f^*\mathcal{F}, (1_G \times f)^*\alpha)$
给出，它定义一个从拟凝聚 $G$-等变层（在 $Y$ 上）的范畴到拟凝聚
$G$-等变层（在 $X$ 上）的范畴的函子。
\end{lemma}

\begin{proof}
略。
\end{proof}
\section{代数空间中的群胚}
\label{spaces-groupoids-section-groupoids}

\noindent
记号参见《群胚》第 \ref{groupoids-section-groupoids} 节。

\begin{definition}
\label{spaces-groupoids-definition-groupoid}
设 $B \to S$ 如第 \ref{spaces-groupoids-section-notation} 节所示。
\begin{enumerate}
\item 一个\emph{$B$ 上代数空间中的群胚}，是一个五元组
$(U, R, s, t, c)$，其中 $U$ 和 $R$ 是 $B$ 上的代数空间，而
$s, t : R \to U$ 和 $c : R \times_{s, U, t} R \to R$ 是 $B$ 上
代数空间的态射，并具有如下性质：对每个概形 $T$（在 $B$ 上），五元组
$$
(U(T), R(T), s, t, c)
$$
是群胚范畴。
\item 代数空间中的群胚的一个态射
$f : (U, R, s, t, c) \to (U', R', s', t', c')$
（取在 $B$ 上），由代数空间的态射 $f : U \to U'$ 和
$f : R \to R'$（在 $B$ 上）给出，并具有如下性质：对每个概形
$T$（在 $B$ 上），这些映射 $f$ 定义一个从群胚范畴
$(U(T), R(T), s, t, c)$ 到群胚范畴
$(U'(T), R'(T), s', t', c')$ 的函子。
\end{enumerate}
\end{definition}

\noindent
设 $(U, R, s, t, c)$ 为代数空间中的群胚（在 $B$ 上）。注意，存在唯一的
代数空间态射 $e : U \to R$ 和 $i : R \to R$（在 $B$ 上），使得对
每个概形 $T$（在 $B$ 上），诱导映射
% P09-ERR-0880
$e : U(T) \to R(T)$ 是恒等，而
% P09-ERR-0881
$i : R(T) \to R(T)$ 是该群胚范畴的逆。七元组
$(U, R, s, t, c, e, i)$ 满足若干交换图表，它们分别对应于《群胚》第
\ref{groupoids-section-groupoids} 节的公理 (1)、(2)(a)、(2)(b)、(3)(a)
和 (3)(b)。反过来，给定具有这些性质的七元组，五元组
$(U, R, s, t, c)$ 就是代数空间中的群胚（在 $B$ 上）。注意，$i$ 是
同构，而 $e$ 同时是 $s$ 和 $t$ 的截面。此外，对一个 $B$ 上代数空间
中的群胚，我们记
$$
j = (t, s) : R \longrightarrow U \times_B U
$$
这与上面第 \ref{spaces-groupoids-section-equivalence-relations} 节的约定相容。
我们有时说“设 $(U, R, s, t, c, e, i)$ 为代数空间中的群胚（在 $B$
上）”，以强调恒等态射和逆态射的存在。

\begin{lemma}
\label{spaces-groupoids-lemma-groupoid-pre-equivalence}
设 $B \to S$ 如第 \ref{spaces-groupoids-section-notation} 节所示。给定代数空间中的
群胚 $(U, R, s, t, c)$（在 $B$ 上），态射
$j : R \to U \times_B U$ 是预等价关系。
\end{lemma}

\begin{proof}
略。这是一个很好的定义练习。
\end{proof}

\begin{lemma}
\label{spaces-groupoids-lemma-equivalence-groupoid}
设 $B \to S$ 如第 \ref{spaces-groupoids-section-notation} 节所示。给定等价关系
$j : R \to U \times_B U$（在 $B$ 上），存在唯一的方式将它扩充成
代数空间中的群胚 $(U, R, s, t, c)$（在 $B$ 上）。
\end{lemma}

\begin{proof}
略。这是一个很好的定义练习。
\end{proof}

\begin{lemma}
\label{spaces-groupoids-lemma-diagram}
设 $B \to S$ 如第 \ref{spaces-groupoids-section-notation} 节所示。设
$(U, R, s, t, c)$ 为代数空间中的群胚（在 $B$ 上）。在交换图表
$$
\xymatrix{
& U & \\
R \ar[d]_s \ar[ru]^t &
R \times_{s, U, t} R
\ar[l]^-{\text{pr}_0} \ar[d]^{\text{pr}_1} \ar[r]_-c &
R \ar[d]^s \ar[lu]_t \\
U & R \ar[l]_t \ar[r]^s & U
}
$$
中，下方两个方块都是纤维积方块。此外，上方的三角形（实际上是方块）也是
笛卡尔的。
\end{lemma}

\begin{proof}
略。这是一个关于定义以及代数几何中函子点观点的练习。
\end{proof}

\begin{lemma}
\label{spaces-groupoids-lemma-diagram-pull}
设 $B \to S$ 如第 \ref{spaces-groupoids-section-notation} 节所示。设
$(U, R, s, t, c, e, i)$ 为代数空间中的群胚（在 $B$ 上）。图表
\begin{equation}
\label{spaces-groupoids-equation-pull}
\xymatrix{
R \times_{t, U, t} R
\ar@<1ex>[r]^-{\text{pr}_1} \ar@<-1ex>[r]_-{\text{pr}_0}
\ar[d]_{\text{pr}_0 \times c \circ (i, 1)} &
R \ar[r]^t \ar[d]^{\text{id}_R} &
U \ar[d]^{\text{id}_U} \\
R \times_{s, U, t} R
\ar@<1ex>[r]^-c \ar@<-1ex>[r]_-{\text{pr}_0} \ar[d]_{\text{pr}_1} &
R \ar[r]^t \ar[d]^s &
U \\
R \ar@<1ex>[r]^s \ar@<-1ex>[r]_t &
U
}
\end{equation}
交换。上方两行通过给出的竖直映射彼此同构。左下方两个方块是笛卡尔的。
\end{lemma}

\begin{proof}
图表的交换性由群胚的公理得到。注意，用群胚的语言说，左上方的竖直箭头把
具有相同目标的一对态射 $(\alpha, \beta)$ 映到一对态射
$(\alpha, \alpha^{-1} \circ \beta)$。在任意群胚中，这定义了
$$
\text{Arrows} \times_{t, \text{Ob}, t} \text{Arrows}
$$
与
$$
\text{Arrows} \times_{s, \text{Ob}, t} \text{Arrows}
$$
之间的双射。因此得到引理的第二个断言。最后一个断言由引理
\ref{spaces-groupoids-lemma-diagram} 得到。
\end{proof}

\begin{lemma}
\label{spaces-groupoids-lemma-base-change-groupoid}
设 $B \to S$ 如第 \ref{spaces-groupoids-section-notation} 节所示。设
$(U, R, s, t, c)$ 为代数空间中的群胚（在 $B$ 上）。设
$B' \to B$ 为代数空间的态射。那么基变换 $U' = B' \times_B U$、
$R' = B' \times_B R$，连同基变换 $s'$、$t'$、$c'$（分别来自态射
$s, t, c$），构成代数空间中的群胚 $(U', R', s', t', c')$（在 $B'$ 上）；
并且这些投影确定代数空间中的群胚的一个态射
$$
(U', R', s', t', c') \to (U, R, s, t, c)
$$
（取在 $B$ 上）。
\end{lemma}

\begin{proof}
略。提示：
$R' \times_{s', U', t'} R' = B' \times_B (R \times_{s, U, t} R)$。
\end{proof}
\section{群胚上的拟凝聚层}
\label{spaces-groupoids-section-groupoids-quasi-coherent}

\noindent
请与《群胚》第 \ref{groupoids-section-groupoids-quasi-coherent} 节比较。

\begin{definition}
\label{spaces-groupoids-definition-groupoid-module}
设 $B \to S$ 如第 \ref{spaces-groupoids-section-notation} 节所示。设
$(U, R, s, t, c)$ 为代数空间中的群胚（在 $B$ 上）。一个
\emph{群胚 $(U, R, s, t, c)$ 上的拟凝聚模}，是一个偶
$(\mathcal{F}, \alpha)$，其中 $\mathcal{F}$ 是拟凝聚
$\mathcal{O}_U$-模，而 $\alpha$ 是 $\mathcal{O}_R$-模映射
$$
\alpha : t^*\mathcal{F} \longrightarrow s^*\mathcal{F}
$$
并且满足：
\begin{enumerate}
\item 图表
$$
\xymatrix{
& \text{pr}_1^*t^*\mathcal{F} \ar[r]_-{\text{pr}_1^*\alpha} &
\text{pr}_1^*s^*\mathcal{F} \ar@{=}[rd] & \\
\text{pr}_0^*s^*\mathcal{F} \ar@{=}[ru] & & & c^*s^*\mathcal{F} \\
& \text{pr}_0^*t^*\mathcal{F} \ar[lu]^{\text{pr}_0^*\alpha} \ar@{=}[r] &
c^*t^*\mathcal{F} \ar[ru]_{c^*\alpha}
}
$$
% P09-ERR-0882
在 $\mathcal{O}_{R \times_{s, U, t} R}$-模范畴中交换；
\item 拉回
$$
e^*\alpha : \mathcal{F} \longrightarrow \mathcal{F}
$$
是恒等映射。
\end{enumerate}
请与引理 \ref{spaces-groupoids-lemma-diagram} 的交换图表比较。
\end{definition}

\noindent
第一个图表的交换性迫使算子 $e^*\alpha$ 为幂等。因此第二个条件可以改述为
$e^*\alpha$ 是同构。事实上，该条件蕴含 $\alpha$ 是同构。

\begin{lemma}
\label{spaces-groupoids-lemma-isomorphism}
设 $B \to S$ 如第 \ref{spaces-groupoids-section-notation} 节所示。设
$(U, R, s, t, c)$ 为代数空间中的群胚（在 $B$ 上）。如果
$(\mathcal{F}, \alpha)$ 是群胚 $(U, R, s, t, c)$ 上的拟凝聚模，
那么 $\alpha$ 是同构。
\end{lemma}

\begin{proof}
把定义 \ref{spaces-groupoids-definition-groupoid-module} 中的交换图表沿态射
$(i, 1) : R \to R \times_{s, U, t} R$ 拉回。于是可见
$i^*\alpha \circ \alpha = s^*e^*\alpha$。沿态射 $(1, i)$ 拉回，
得到关系
$\alpha \circ i^*\alpha = t^*e^*\alpha$。由第二个假设，这些态射
是恒等态射。因此 $i^*\alpha$ 是 $\alpha$ 的逆。
\end{proof}

\begin{lemma}
\label{spaces-groupoids-lemma-pullback}
设 $B \to S$ 如第 \ref{spaces-groupoids-section-notation} 节所示。考虑一个态射
$f : (U, R, s, t, c) \to (U', R', s', t', c')$，
% P09-ERR-0883
它是代数空间中的群胚的态射（在 $B$ 上）。
% P09-ERR-0884
拉回 $f^*$ 由
$$
(\mathcal{F}, \alpha) \mapsto (f^*\mathcal{F}, f^*\alpha)
$$
给出，并定义一个从群胚 $(U', R', s', t', c')$ 上的拟凝聚层范畴到群胚
$(U, R, s, t, c)$ 上的拟凝聚层范畴的函子。
\end{lemma}

\begin{proof}
略。
\end{proof}

\begin{lemma}
\label{spaces-groupoids-lemma-pushforward}
设 $B \to S$ 如第 \ref{spaces-groupoids-section-notation} 节所示。考虑代数空间中的群胚的
一个态射
$f : (U, R, s, t, c) \to (U', R', s', t', c')$（在 $B$ 上）。
假设：
\begin{enumerate}
\item $f : U \to U'$ 拟紧且拟分离；
\item 方块
$$
\xymatrix{
R \ar[d]_t \ar[r]_f & R' \ar[d]^{t'} \\
U \ar[r]^f & U'
}
$$
是笛卡尔的；
\item $s'$ 和 $t'$ 平坦。
\end{enumerate}
% P09-ERR-0885
前推 $f_*$ 由
$$
(\mathcal{F}, \alpha) \mapsto (f_*\mathcal{F}, f_*\alpha)
$$
给出，并定义一个从群胚 $(U, R, s, t, c)$ 上的拟凝聚层范畴到群胚
$(U', R', s', t', c')$ 上的拟凝聚层范畴的函子；这个函子是引理
\ref{spaces-groupoids-lemma-pullback} 所定义拉回的右伴随。
\end{lemma}

\begin{proof}
由于 $U \to U'$ 拟紧且拟分离，可知 $f_*$ 把拟凝聚层变为拟凝聚层
（《空间的态射》引理 \ref{spaces-morphisms-lemma-pushforward}）。
此外，由于方块
$$
\vcenter{
\xymatrix{
R \ar[d]_t \ar[r]_f & R' \ar[d]^{t'} \\
U \ar[r]^f & U'
}
}
\quad\text{和}\quad
\vcenter{
\xymatrix{
R \ar[d]_s \ar[r]_f & R' \ar[d]^{s'} \\
U \ar[r]^f & U'
}
}
$$
是笛卡尔的，得到
$(t')^*f_*\mathcal{F} = f_*t^*\mathcal{F}$ 和
$(s')^*f_*\mathcal{F} = f_*s^*\mathcal{F}$，见《空间的上同调》引理
\ref{spaces-cohomology-lemma-flat-base-change-cohomology}。因此，把
$f_*\alpha$ 看作映射
$(t')^*f_*\mathcal{F} \to (s')^*f_*\mathcal{F}$ 是有意义的。
类似的论证表明 $f_*\alpha$ 满足余圈条件。这个函子与拉回函子互为伴随，
因为环化空间上模的拉回与前推互为伴随。略去一些细节。
\end{proof}

\begin{lemma}
\label{spaces-groupoids-lemma-colimits}
设 $B \to S$ 如第 \ref{spaces-groupoids-section-notation} 节所示。设
$(U, R, s, t, c)$ 为代数空间中的群胚（在 $B$ 上）。群胚
$(U, R, s, t, c)$ 上的拟凝聚模范畴有余极限。
\end{lemma}

\begin{proof}
设 $i \mapsto (\mathcal{F}_i, \alpha_i)$ 为以范畴 $\mathcal{I}$ 为
指标的图表。可以形成余极限
$\mathcal{F} = \colim \mathcal{F}_i$；它是 $U$ 上的拟凝聚层，见
《空间的性质》引理
\ref{spaces-properties-lemma-properties-quasi-coherent}。由于余极限与
拉回交换，可知
$s^*\mathcal{F} = \colim s^*\mathcal{F}_i$，类似地
$t^*\mathcal{F} = \colim t^*\mathcal{F}_i$。因此可以置
$\alpha = \colim \alpha_i$。我们略去如下断言的证明：
$(\mathcal{F}, \alpha)$ 是拟凝聚模范畴中该图表的余极限，该范畴取在群胚
$(U, R, s, t, c)$ 上。
\end{proof}

\begin{lemma}
\label{spaces-groupoids-lemma-abelian}
设 $B \to S$ 如第 \ref{spaces-groupoids-section-notation} 节所示。设
$(U, R, s, t, c)$ 为代数空间中的群胚（在 $B$ 上）。如果 $s$、$t$
平坦，那么群胚 $(U, R, s, t, c)$ 上的拟凝聚模范畴是 Abel 范畴。
\end{lemma}

\begin{proof}
设 $\varphi : (\mathcal{F}, \alpha) \to (\mathcal{G}, \beta)$ 是
群胚 $(U, R, s, t, c)$ 上拟凝聚模的同态。由于 $s$ 平坦，可知
$$
0 \to s^*\Ker(\varphi)
\to s^*\mathcal{F} \to s^*\mathcal{G} \to s^*\Coker(\varphi) \to 0
$$
正合；沿 $t$ 拉回时也一样。因此 $\alpha$ 和 $\beta$ 诱导同构
$\kappa : t^*\Ker(\varphi) \to s^*\Ker(\varphi)$ 和
$\lambda : t^*\Coker(\varphi) \to s^*\Coker(\varphi)$，它们满足
余圈条件。于是直接验证可知 $(\Ker(\varphi), \kappa)$ 和
$(\Coker(\varphi), \lambda)$ 分别是拟凝聚模范畴中的核和余核，该范畴
取在群胚 $(U, R, s, t, c)$ 上。此外，由于等式
$\Coim(\varphi) = \Im(\varphi)$ 在 $U$ 上成立，所以这里也成立。
\end{proof}
\section{拟凝聚模的余极限}
\label{spaces-groupoids-section-colimits}

\noindent
本节对应于《群胚》第 \ref{groupoids-section-colimits} 节。

\begin{lemma}
\label{spaces-groupoids-lemma-construct-quasi-coherent}
设 $B \to S$ 如第 \ref{spaces-groupoids-section-notation} 节所示。设
$(U, R, s, t, c)$ 为代数空间中的群胚（在 $B$ 上）。假设 $s, t$ 平坦、
拟紧且拟分离。对任意拟凝聚模 $\mathcal{G}$（在 $U$ 上），存在典范同构
$\alpha : t^*s_*t^*\mathcal{G} \to s^*s_*t^*\mathcal{G}$，它使
$(s_*t^*\mathcal{G}, \alpha)$ 成为群胚 $(U, R, s, t, c)$ 上的
拟凝聚模。这个构造定义函子
$$
\QCoh(\mathcal{O}_U) \longrightarrow \QCoh(U, R, s, t, c)
$$
它是遗忘函子 $(\mathcal{F}, \beta) \mapsto \mathcal{F}$ 的右伴随。
\end{lemma}

\begin{proof}
沿拟紧拟分离态射前推拟凝聚模，所得仍为拟凝聚模，见《空间的态射》引理
\ref{spaces-morphisms-lemma-pushforward}。因此 $s_*t^*\mathcal{G}$ 是
拟凝聚的。使用引理 \ref{spaces-groupoids-lemma-diagram} 的记号，有
$$
t^*s_*t^*\mathcal{G} =
\text{pr}_{1, *}\text{pr}_0^*t^*\mathcal{G} =
\text{pr}_{1, *}c^*t^*\mathcal{G} =
s^*s_*t^*\mathcal{G}
$$
中间的等号来自等式
$t \circ c = t \circ \text{pr}_0$，这是态射
$R \times_{s, U, t} R \to U$ 的等式；第一个和最后一个等号来自这些步骤中
基变换与前推交换，见《空间的上同调》引理
\ref{spaces-cohomology-lemma-flat-base-change-cohomology}。

\medskip\noindent
为了验证定义 \ref{spaces-groupoids-definition-groupoid-module} 中 $\alpha$ 的余圈条件以及
伴随性质，我们用另一种方式描述构造
$\mathcal{G} \mapsto (s_*t^*\mathcal{G}, \alpha)$。考虑
% P09-ERR-0886
代数空间中的群胚
% P09-ERR-0888
$(R, R \times_{t, U, t} R, \text{pr}_0, \text{pr}_1, \text{pr}_{02})$；
它来自等价关系 $R \times_{t, U, t} R$（定义在 $R$ 上），见引理
\ref{spaces-groupoids-lemma-equivalence-groupoid}。存在态射
$$
f :
(R, R \times_{t, U, t} R, \text{pr}_1, \text{pr}_0, \text{pr}_{02})
\longrightarrow
(U, R, s, t, c)
$$
% P09-ERR-0887
这是代数空间中的群胚的态射，由 $s : R \to U$ 和映射
$R \times_{t, U, t} R \to R$ 给出，后一个映射为
$(r_0, r_1) \mapsto r_0^{-1} \circ r_1$；我们略去所需图表交换性的验证。
由于 $t, s : R \to U$ 拟紧、拟分离且平坦，并且有笛卡尔方块
$$
\xymatrix{
R \times_{t, U, t} R \ar[d]_{\text{pr}_0}
\ar[rr]_-{(r_0, r_1) \mapsto r_0^{-1} \circ r_1} & & R \ar[d]^t \\
R \ar[rr]^s & & U
}
$$
由引理 \ref{spaces-groupoids-lemma-diagram-pull} 可知，引理 \ref{spaces-groupoids-lemma-pushforward}
适用于 $f$。因此，沿 $f$ 前推和拉回拟凝聚模是互为伴随的函子。为了完成证明，
我们把这些函子与上面描述的函子等同起来。为此注意，函子
$$
t^* :
\QCoh(\mathcal{O}_U)
\longrightarrow
\QCoh(R, R \times_{t, U, t} R, \text{pr}_1, \text{pr}_0, \text{pr}_{02})
$$
是等价，因为根据拟凝聚层的下降理论，
$\{t : R \to U\}$ 是 fpqc 覆盖；见《空间上的下降》命题
\ref{spaces-descent-proposition-fpqc-descent-quasi-coherent}。

\medskip\noindent
沿 $f$ 前推并预合成等价 $t^*$，把 $\mathcal{G}$ 映到
$(s_*t^*\mathcal{G}, \alpha)$；我们略去如下验证：以这种方式得到的同构
$\alpha$ 与上面构造的同构相同。

\medskip\noindent
沿 $f$ 拉回并后合成等价 $t^*$ 的逆，把 $(\mathcal{F}, \beta)$ 映到
相对于 $\{t : R \to U\}$ 的下降；被下降的模为 $s^*\mathcal{F}$，配备
下降资料 $\gamma$；该资料定义在 $R \times_{t, U, t} R$ 上，并且是
$\beta$ 沿映射
$(r_0, r_1) \mapsto r_0^{-1} \circ r_1$ 的拉回。考虑同构
$\beta : t^*\mathcal{F} \to s^*\mathcal{F}$。$t^*\mathcal{F}$ 上相对于
$\{t : R \to U\}$ 的典范下降资料（《空间上的下降》
定义 \ref{spaces-descent-definition-descent-datum-effective-quasi-coherent}）
经由 $\beta$ 变成映射
$$
\text{pr}_0^*s^*\mathcal{F}
\xrightarrow{\text{pr}_0^*\beta^{-1}}
\text{pr}_0^*t^*\mathcal{F}
\xrightarrow{can}
\text{pr}_1^*t^*\mathcal{F}
\xrightarrow{\text{pr}_1^*\beta}
\text{pr}_1^*s^*\mathcal{F}
$$
由于 $\beta$ 满足余圈条件，这等于 $\beta$ 沿映射
$(r_0, r_1) \mapsto r_0^{-1} \circ r_1$ 的拉回。为说明这一点，取定义
\ref{spaces-groupoids-definition-groupoid-module} 中实际的余圈关系，并沿态射
$(\text{pr}_0, c \circ (i, 1)) : R \times_{t, U, t} R \to R \times_{s, U, t} R$
将它拉回；该态射也出现在引理 \ref{spaces-groupoids-lemma-diagram-pull} 的交换图表中。
由此可知 $(s^*\mathcal{F}, \gamma)$ 同构于
$(t^*\mathcal{F}, can)$。综上，沿 $f$ 拉回并后合成等价 $t^*$ 的逆所得的
函子，同构于遗忘函子 $(\mathcal{F}, \beta) \mapsto \mathcal{F}$。
\end{proof}

\begin{remark}
\label{spaces-groupoids-remark-adjunction-map}
在引理 \ref{spaces-groupoids-lemma-construct-quasi-coherent} 的情形下，记
$$
F : \QCoh(U, R, s, t, c) \to \QCoh(\mathcal{O}_U),\quad
(\mathcal{F}, \beta) \mapsto \mathcal{F}
$$
为遗忘函子，并记
$$
G : \QCoh(\mathcal{O}_U) \to \QCoh(U, R, s, t, c),\quad
\mathcal{G} \mapsto (s_*t^*\mathcal{G}, \alpha)
$$
为引理中构造的右伴随。那么，该伴随的单位
$\eta : \text{id} \to G \circ F$ 在 $(\mathcal{F}, \beta)$ 上的取值由
映射
$$
\mathcal{F} \to s_*s^*\mathcal{F} \xrightarrow{\beta^{-1}} s_*t^*\mathcal{F}
$$
给出。我们略去验证。
\end{remark}

\begin{lemma}
\label{spaces-groupoids-lemma-push-pull}
设 $S$ 为概形。设 $f : Y \to X$ 为 $S$ 上代数空间的态射。设
$\mathcal{F}$ 为拟凝聚 $\mathcal{O}_X$-模，设 $\mathcal{G}$ 为拟凝聚
$\mathcal{O}_Y$-模，并设
$\varphi : \mathcal{G} \to f^*\mathcal{F}$ 为模映射。假设：
\begin{enumerate}
\item $\varphi$ 是单射；
\item $f$ 拟紧、拟分离、平坦且满；
\item $X$、$Y$ 局部 Noether；
\item $\mathcal{G}$ 是凝聚 $\mathcal{O}_Y$-模。
\end{enumerate}
那么 $\mathcal{F} \cap f_*\mathcal{G}$ 由拉回
$$
\xymatrix{
\mathcal{F} \ar[r] & f_*f^*\mathcal{F} \\
\mathcal{F} \cap f_*\mathcal{G} \ar[u] \ar[r] &
f_*\mathcal{G} \ar[u]
}
$$
定义，并且它是凝聚 $\mathcal{O}_X$-模。
\end{lemma}

\begin{proof}
我们将自由使用《空间的上同调》引理
\ref{spaces-cohomology-lemma-coherent-Noetherian} 对凝聚模的刻画，以及
凝聚模构成 $\QCoh(\mathcal{O}_X)$ 的 Serre 子范畴这一事实；见
《空间的上同调》引理
\ref{spaces-cohomology-lemma-coherent-Noetherian-quasi-coherent-sub-quotient}。
如果 $f$ 有截面 $\sigma$，那么可见 $\mathcal{F} \cap f_*\mathcal{G}$
包含在
$\sigma^*\mathcal{G} \to \sigma^*f^*\mathcal{F} = \mathcal{F}$
的像中，因而是凝聚的。一般地，为证明
$\mathcal{F} \cap f_*\mathcal{G}$ 凝聚，
% P09-ERR-0889
只需证明 $f^*(\mathcal{F} \cap f_*\mathcal{G})$ 凝聚（见
《空间上的下降》引理 \ref{spaces-descent-lemma-finite-type-descends}）。
由于 $f$ 平坦，这等于
$f^*\mathcal{F} \cap f^*f_*\mathcal{G}$。由于 $f$ 平坦、拟紧且拟分离，
可知 $f^*f_*\mathcal{G} = p_*q^*\mathcal{G}$，其中
$p, q : Y \times_X Y \to Y$ 是投影；见《空间的上同调》引理
\ref{spaces-cohomology-lemma-flat-base-change-cohomology}。由于 $p$
有截面，结论成立。
\end{proof}

\noindent
设 $B \to S$ 如第 \ref{spaces-groupoids-section-notation} 节所示。设
$(U, R, s, t, c)$ 为代数空间中的群胚（在 $B$ 上）。假设 $U$ 局部
Noether。在下面的引理中，如果拟凝聚层 $(\mathcal{F}, \alpha)$（定义在群胚
$(U, R, s, t, c)$ 上）满足 $\mathcal{F}$ 是凝聚
$\mathcal{O}_U$-模，就称它\emph{凝聚}。

\begin{lemma}
\label{spaces-groupoids-lemma-colimit-coherent}
设 $B \to S$ 如第 \ref{spaces-groupoids-section-notation} 节所示。设
$(U, R, s, t, c)$ 为代数空间中的群胚（在 $B$ 上）。假设：
\begin{enumerate}
\item $U$、$R$ 是 Noether 的；
\item $s, t$ 平坦、拟紧且拟分离。
\end{enumerate}
那么每个拟凝聚模 $(\mathcal{F}, \alpha)$（定义在群胚
$(U, R, s, t, c)$ 上）都是凝聚模的滤过余极限。
\end{lemma}

\begin{proof}
我们将使用《空间的上同调》引理
\ref{spaces-cohomology-lemma-coherent-Noetherian} 对局部 Noether
代数空间上凝聚模的刻画，而不再提及。可以把
$\mathcal{F} = \colim \mathcal{H}_i$ 写成凝聚子模
$\mathcal{H}_i \subset \mathcal{F}$ 的滤过余极限，见《空间的上同调》
引理 \ref{spaces-cohomology-lemma-directed-colimit-coherent}。给定拟凝聚层
$\mathcal{H}$（在 $U$ 上），记
$(s_*t^*\mathcal{H}, \alpha)$ 为引理
\ref{spaces-groupoids-lemma-construct-quasi-coherent} 所给出的群胚
$(U, R, s, t, c)$ 上的拟凝聚层。考虑伴随映射
$(\mathcal{F}, \beta) \to (s_*t^*\mathcal{F}, \alpha)$（在
$\QCoh(U, R, s, t, c)$ 中），见注
\ref{spaces-groupoids-remark-adjunction-map}。置
$$
(\mathcal{F}_i, \beta_i) =
(\mathcal{F}, \beta)
\times_{(s_*t^*\mathcal{F}, \alpha)}
(s_*t^*\mathcal{H}_i, \alpha)
$$
于 $\QCoh(U, R, s, t, c)$ 中。由引理 \ref{spaces-groupoids-lemma-abelian} 的证明，
限制到 $U$ 在 $\QCoh(U, R, s, t, c)$ 上是正合函子，因此得到拉回图表
$$
\xymatrix{
\mathcal{F} \ar[r] & s_*t^*\mathcal{F} \\
\mathcal{F}_i \ar[r] \ar[u] & s_*t^*\mathcal{H}_i \ar[u]
}
$$
换言之，$\mathcal{F}_i = \mathcal{F} \cap s_*t^*\mathcal{H}_i$。
根据注 \ref{spaces-groupoids-remark-adjunction-map} 中伴随映射的描述，该图表同构于图表
$$
\xymatrix{
\mathcal{F} \ar[r] & s_*s^*\mathcal{F} \\
\mathcal{F}_i \ar[r] \ar[u] & s_*t^*\mathcal{H}_i \ar[u]
}
$$
其中右侧竖直箭头是把 $s_*$ 施于映射
$$
t^*\mathcal{H}_i \to t^*\mathcal{F} \xrightarrow{\beta} s^*\mathcal{F}
$$
所得的结果。该箭头是单射，因为 $t$ 是平坦态射。由引理
\ref{spaces-groupoids-lemma-push-pull} 可知 $\mathcal{F}_i$ 凝聚。最后，因为 $s$ 拟紧且
拟分离，可知 $s_*$ 与余极限交换（见《概形的上同调》引理
\ref{coherent-lemma-colimit-cohomology}）。因此
$s_*t^*\mathcal{F} = \colim s_*t^*\mathcal{H}_i$，从而如期有
$(\mathcal{F}, \beta) = \colim (\mathcal{F}_i, \beta_i)$。
\end{proof}

\section{拟凝聚层中的晶体}
\label{spaces-groupoids-section-crystals}

\noindent
设 $(I, \Phi, j)$ 由集合 $I$ 和预关系
$j : \Phi \to I \times I$ 组成。假设对每个 $i \in I$ 给定概形 $X_i$，
并对每个 $\phi \in \Phi$ 给定概形态射
$f_\phi : X_{i'} \to X_i$，其中 $j(\phi) = (i, i')$。置
$X = (\{X_i\}_{i \in I}, \{f_\phi\}_{\phi \in \Phi})$。
所谓 $X$ 上拟凝聚模的\emph{晶体}，是指如下规则：它对每个
% P09-ERR-0890
$i \in \Ob(\mathcal{I})$ 指定拟凝聚层 $\mathcal{F}_i$（在 $X_i$ 上），
并对每个 $\phi \in \Phi$（满足 $j(\phi) = (i, i')$）指定拟凝聚层同构
$$
\alpha_\phi : f_\phi^*\mathcal{F}_i \longrightarrow \mathcal{F}_{i'}
$$
（定义在 $X_{i'}$ 上）。
这些拟凝聚模晶体构成加性范畴 $\textit{CQC}(X)$\footnote{可以选出一族三元组
$\phi, \phi', \phi'' \in \Phi$，它们满足
$j(\phi) = (i, i')$、$j(\phi') = (i', i'')$ 和
$j(\phi'') = (i, i'')$，并且
$f_{\phi''} = f_\phi \circ f_{\phi'}$；然后要求这些三元组满足
$\alpha_{\phi'} \circ f_{\phi'}^*\alpha_\phi = \alpha_{\phi''}$。
这会定义一个加性子范畴。例如，资料 $(I, \Phi)$ 可以是某个指标范畴的
对象集与箭头集，而 $X$ 可以是该指标范畴上的概形图表。引理
\ref{spaces-groupoids-lemma-crystals-in-quasi-coherent-modules} 的结论立即给出相应子范畴中的结论。}。
这个范畴有余极限（证明与引理 \ref{spaces-groupoids-lemma-colimits} 的证明相同）。
如果所有态射 $f_\phi$ 都平坦，那么 $\textit{CQC}(X)$ 是 Abel 范畴
（证明与引理 \ref{spaces-groupoids-lemma-abelian} 的证明相同）。设 $\kappa$ 为基数。
称拟凝聚模晶体 $\mathcal{F}$（在 $X$ 上）是\emph{$\kappa$-生成的}，
如果它的每个 $\mathcal{F}_i$ 都是 $\kappa$-生成的（见《性质》定义
\ref{properties-definition-kappa-generated}）。

\begin{lemma}
\label{spaces-groupoids-lemma-crystals-in-quasi-coherent-modules}
在上述情形中，如果所有态射 $f_\phi$ 都平坦，那么存在基数 $\kappa$，
使得每个对象
$(\{\mathcal{F}_i\}_{i \in I}, \{\alpha_\phi\}_{\phi \in \Phi})$
（属于 $\textit{CQC}(X)$）
都是其 $\kappa$-生成子模的有向余极限。
\end{lemma}

\begin{proof}
在本引理及其证明中，
$(\{\mathcal{F}_i\}_{i \in I}, \{\alpha_\phi\}_{\phi \in \Phi})$
的\emph{子模}，是指给定拟凝聚子模
$\mathcal{G}_i \subset \mathcal{F}_i$（对所有 $i$），使得
$\alpha_\phi(f_\phi^*\mathcal{G}_i) = \mathcal{G}_{i'}$
作为 $\mathcal{F}_{i'}$ 的子层，并且这对所有 $\phi \in \Phi$ 成立。
这个定义有意义，因为 $f_\phi$ 平坦时，
拉回 $f^*_\phi$ 是正合的，也就是说它保持子层。
% P09-ERR-0891
本证明是《性质》引理 \ref{properties-lemma-colimit-kappa} 的证明的一个变体。
建议读者先阅读该证明。

\medskip\noindent
我们声称，只需在所有概形 $X_i$ 都仿射的情形下证明本引理。为说明这一点，令
$$
J = \coprod\nolimits_{i \in I} \{U \subset X_i\text{ 为仿射开集}\}
$$
并令
% P09-ERR-0892
\begin{align*}
\Psi = & \coprod\nolimits_{\phi \in \Phi}
\{
(U, V) \mid
U \subset X_i, V \subset X_{i'}\text{ 为仿射开集且 } f_\phi(U) \subset V
\} \\
&
\amalg \coprod\nolimits_{i \in I}
\{
(U, U') \mid
U, U' \subset X_i\text{ 为仿射开集且 } U \subset U'
\}.
\end{align*}
给它配备显然的映射 $\Psi \to J \times J$。于是我们的
$(\mathcal{F}, \alpha)$ 诱导拟凝聚层晶体
$(\{\mathcal{H}_j\}_{j \in J}, \{\beta_\psi\}_{\psi \in \Psi})$
（在 $Y = (J, \Psi)$ 上）：置
$\mathcal{H}_{(i, U)} = \mathcal{F}_i|_U$（对 $(i, U) \in J$），并置
$\beta_\psi$（对 $\psi \in \Psi$）如下：
% P09-ERR-0893
它等于 $\alpha_\phi$ 在 $U$ 上的限制（若
$\psi = (\phi, U, V)$），并等于
$\text{id} : (\mathcal{F}_i|_{U'})|_U \to \mathcal{F}_i|_U$
（当 $\psi = (i, U, U')$）。此外，
$(\{\mathcal{H}_j\}_{j \in J}, \{\beta_\psi\}_{\psi \in \Psi})$
的子模按 $1$ 对 $1$ 地对应于
$(\{\mathcal{F}_i\}_{i \in I}, \{\alpha_\phi\}_{\phi \in \Phi})$
的子模。我们略去证明（提示：使用《层》第
\ref{sheaves-section-bases} 节）。此外，显然，如果 $\kappa$ 对 $Y$ 可行，
那么同一个 $\kappa$ 对 $X$ 也可行（由 $\kappa$-生成模的定义）。
% P09-ERR-0894
因此，只需对 $Y$ 上拟凝聚层晶体证明本引理。

\medskip\noindent
假设所有概形 $X_i$ 都仿射。取无限基数 $\kappa$，使其大于 $I$ 或
$\Phi$ 的基数。设
$(\{\mathcal{F}_i\}_{i \in I}, \{\alpha_\phi\}_{\phi \in \Phi})$
为 $\textit{CQC}(X)$ 的对象。对每个 $i$ 写成
$X_i = \Spec(A_i)$ 和 $M_i = \Gamma(X_i, \mathcal{F}_i)$。
对每个 $\phi \in \Phi$（满足 $j(\phi) = (i, i')$），映射
$\alpha_\phi$ 转化为 $A_{i'}$-模同构
$$
\alpha_\phi : M_i \otimes_{A_i} A_{i'} \longrightarrow M_{i'}
$$
利用选择公理，选取规则
% P09-ERR-0895
$$
(\phi, m) \longmapsto S(\phi, m')
$$
其中源是满足如下条件的偶 $(\phi, m')$ 的集合：
$\phi \in \Phi$、$j(\phi) = (i, i')$ 且 $m' \in M_{i'}$；其输出是
有限子集 $S(\phi, m') \subset M_i$，使得
$$
m' = \alpha_\phi\left(\sum\nolimits_{m \in S(\phi, m')} m \otimes a'_m\right)
$$
对某些 $a'_m \in A_{i'}$ 成立。

\medskip\noindent
作出这些选择后，我们声称，任意 $\mathcal{F}_i$ 在任意 $X_i$ 上的任意
截面都属于某个 $\kappa$-生成子模。为说明这一点，假设给定子集族
$\mathcal{S} = \{S_i\}_{i \in I}$，其中每个
$S_i \subset M_i$ 的基数至多为 $\kappa$。定义新子集族
$\mathcal{S}' = \{S'_i\}_{i \in I}$，其中
$$
S'_i = S_i \cup
\bigcup\nolimits_{(\phi, m'),\ j(\phi) = (i, i'),\ m' \in S_{i'}} S(\phi, m')
$$
注意每个 $S'_i$ 的基数仍至多为 $\kappa$。置
$\mathcal{S}^{(0)} = \mathcal{S}$、$\mathcal{S}^{(1)} = \mathcal{S}'$，
并递归地置
$\mathcal{S}^{(n + 1)} = (\mathcal{S}^{(n)})'$。再置
$S_i^{(\infty)} = \bigcup_{n \geq 0} S_i^{(n)}$ 和
$\mathcal{S}^{(\infty)} = \{S_i^{(\infty)}\}_{i \in I}$。
根据构造，对每个 $\phi \in \Phi$（满足 $j(\phi) = (i, i')$）及每个
$m' \in S^{(\infty)}_{i'}$，都能把 $m'$ 写成像
$\alpha_\phi(m \otimes 1)$ 的有限线性组合，其中
$m \in S_i^{(\infty)}$。因此，令 $N_i$ 为 $A_i$-子模；它是在
$M_i$ 中由 $S_i^{(\infty)}$ 生成的。则相应的拟凝聚子模
% P09-ERR-0896
$\widetilde{N_i} \subset \mathcal{F}_i$ 构成一个 $\kappa$-生成子模。
这就完成了证明。
\end{proof}

\begin{lemma}
\label{spaces-groupoids-lemma-set-generators}
设 $B \to S$ 如第 \ref{spaces-groupoids-section-notation} 节所示。设
$(U, R, s, t, c)$ 为代数空间中的群胚（在 $B$ 上）。如果 $s$、$t$
平坦，那么存在集合 $T$ 和一族对象
$(\mathcal{F}_t, \alpha_t)_{t \in T}$（属于 $\QCoh(U, R, s, t, c)$），使得每个对象
$(\mathcal{F}, \alpha)$ 都是其如下子模的有向余极限：这些子模同构于
对象 $(\mathcal{F}_t, \alpha_t)$ 中的某一个。
\end{lemma}

\begin{proof}
本引理推广了《群胚》引理 \ref{groupoids-lemma-colimit-kappa}；
后者处理概形中的群胚情形。我们不能直接使用同一个论证，因此使用上面发展的
“拟凝聚层晶体”材料。

\medskip\noindent
选取概形 $W$ 和满 \'etale 态射 $W \to U$。选取概形 $V$ 和满
\'etale 态射 $V \to W \times_{U, s} R$。选取概形 $V'$ 和满
\'etale 态射 $V' \to R \times_{t, U} W$。考虑概形族
$$
I = \{W, W \times_U W, V, V', V \times_R V'\}
$$
以及概形态射集
$$
\Phi = \{\text{pr}_i : W \times_U W \to W, V \to W, V' \to W,
V \times_R V' \to V, V \times_R V' \to V'\}
$$
置 $X = (I, \Phi)$。回忆我们已定义了范畴
$\textit{CQC}(X)$，即 $X$ 上拟凝聚层晶体的范畴。存在函子
$$
\QCoh(U, R, s, t, c) \longrightarrow \textit{CQC}(X)
$$
该函子把 $(\mathcal{F}, \alpha)$ 分别映到：层
$\mathcal{F}|_W$（在 $W$ 上）；层
$\mathcal{F}|_{W \times_U W}$（在 $W \times_U W$ 上）；
$\mathcal{F}$ 经由 $V \to W \times_{U, s} R \to W \to U$
的拉回（在 $V$ 上）；$\mathcal{F}$ 经由
$V' \to R \times_{t, U} W \to W \to U$ 的拉回（在 $V'$ 上）；
最后是 $\mathcal{F}$ 经由
$V \times_R V' \to V \to W \times_{U, s} R \to W \to U$
的拉回（在 $V \times_R V'$ 上）。比较映射
$\{\alpha_\phi\}_{\phi \in \Phi}$ 除一个以外都取显然的映射
（来自拉回的结合性）；对映射
$\phi = \text{pr}_{V'} : V \times_R V' \to V'$，则使用
$\alpha : t^*\mathcal{F} \to s^*\mathcal{F}$ 在
$V \times_R V'$ 上的拉回。下列交换图表说明这个定义有意义：
$$
\xymatrix{
& V \times_R V' \ar[ld] \ar[rd] \\
V \ar[rd] \ar[dd] & & V' \ar[ld] \ar[dd] \\
& R \ar@<-1ex>[dd]_s \ar@<1ex>[dd]^t \\
W \ar[rd] & & W \ar[ld] \\
& U
}
$$
上面显示的函子不是范畴等价。不过，因为 $W \to U$ 是满
\'etale 态射，它是忠实的\footnote{事实上该函子是全忠实的，但我们不需要这一点。}。
由于上图中的所有态射都平坦，可见它是 Abel 范畴间的正合函子。此外，我们声称，
% P09-ERR-0897
给定 $(\mathcal{F}, \alpha)$，其像为
$(\{\mathcal{F}_i\}_{i \in I}, \{\alpha_\phi\}_{\phi \in \Phi})$，
则存在 $1$ 对 $1$ 的对应：$(\mathcal{F}, \alpha)$ 的拟凝聚子模与
$(\{\mathcal{F}_i\}_{i \in I}, \{\alpha_\phi\}_{\phi \in \Phi})$
的子模相对应。具体而言，给定
$(\{\mathcal{F}_i\}_{i \in I}, \{\alpha_\phi\}_{\phi \in \Phi})$
的子模，它在 $W$ 上的部分与投影映射
$W \times_U W \to W$ 相容，这保证该子模来自 $\mathcal{F}$ 的一个
拟凝聚子模（由《空间的性质》命题
\ref{spaces-properties-proposition-quasi-coherent}）；它与
$\alpha_{\text{pr}_{V'}}$ 相容，这保证该子层与 $\alpha$ 相容
（细节略）。

\medskip\noindent
选取基数 $\kappa$，如引理
\ref{spaces-groupoids-lemma-crystals-in-quasi-coherent-modules} 对系统 $X = (I, \Phi)$ 所给。
由《性质》引理 \ref{properties-lemma-set-of-iso-classes} 可直接看出，
$\kappa$-生成的拟凝聚层晶体（在 $X$ 上）的同构类构成一个集合。因此结论显然。
\end{proof}






\section{群胚与群空间}
\label{spaces-groupoids-section-groupoids-group-spaces}

\noindent
请与《群胚》第 \ref{groupoids-section-groupoids-group-schemes} 节比较。

\begin{lemma}
\label{spaces-groupoids-lemma-groupoid-from-action}
设 $B \to S$ 如第 \ref{spaces-groupoids-section-notation} 节所示。设 $(G, m)$ 为
$B$ 上的群代数空间，其单位元为 $e_G$、逆元为 $i_G$。设 $X$ 为
$B$ 上的代数空间，并设 $a : G \times_B X \to X$ 为 $G$ 在
$X$ 上的作用（在 $B$ 上）。那么按如下方式得到代数空间中的群胚
$(U, R, s, t, c, e, i)$（在 $B$ 上）：
\begin{enumerate}
\item 置 $U = X$ 和 $R = G \times_B X$。
\item 置 $s : R \to U$ 等于 $(g, x) \mapsto x$。
\item 置 $t : R \to U$ 等于 $(g, x) \mapsto a(g, x)$。
\item 置 $c : R \times_{s, U, t} R \to R$ 等于
$((g, x), (g', x')) \mapsto (m(g, g'), x')$。
\item 置 $e : U \to R$ 等于 $x \mapsto (e_G(x), x)$。
\item 置 $i : R \to R$ 等于 $(g, x) \mapsto (i_G(g), a(g, x))$。
\end{enumerate}
\end{lemma}

\begin{proof}
略。提示：只需在集合层次证明上述构造可行。为此，使用引理上方的描述：
$g$ 是从 $v$ 到 $a(g, v)$ 的箭头。
\end{proof}

\begin{lemma}
\label{spaces-groupoids-lemma-action-groupoid-modules}
设 $B \to S$ 如第 \ref{spaces-groupoids-section-notation} 节所示。设 $(G, m)$ 为
$B$ 上的群代数空间。设 $X$ 为 $B$ 上的代数空间，并设
$a : G \times_B X \to X$ 为 $G$ 在 $X$ 上的作用（在 $B$ 上）。设
$(U, R, s, t, c)$ 为引理 \ref{spaces-groupoids-lemma-groupoid-from-action} 中构造的
代数空间中的群胚。规则
$(\mathcal{F}, \alpha) \mapsto (\mathcal{F}, \alpha)$ 给出如下两个
范畴之间的等价：$G$-等变 $\mathcal{O}_X$-模的范畴，与群胚
$(U, R, s, t, c)$ 上拟凝聚模的范畴。
\end{lemma}

\begin{proof}
该断言有意义，因为 $t = a$ 且 $s = \text{pr}_1$（作为态射
$R = G \times_B X \to X$）；见定义 \ref{spaces-groupoids-definition-equivariant-module}
和 \ref{spaces-groupoids-definition-groupoid-module}。利用引理
\ref{spaces-groupoids-lemma-groupoid-from-action} 中的转译，两个定义中的交换性要求
恰好相符。
\end{proof}






\section{稳定子群代数空间}
\label{spaces-groupoids-section-stabilizer}

\noindent
请与《群胚》第 \ref{groupoids-section-stabilizer} 节比较。给定代数空间中的群胚，
可如下得到一个群代数空间。

\begin{lemma}
\label{spaces-groupoids-lemma-groupoid-stabilizer}
设 $B \to S$ 如第 \ref{spaces-groupoids-section-notation} 节所示。设
$(U, R, s, t, c)$ 为代数空间中的群胚（在 $B$ 上）。代数空间 $G$ 由笛卡儿方块
$$
\xymatrix{
G \ar[r] \ar[d] & R \ar[d]^{j = (t, s)} \\
U \ar[r]^-{\Delta} & U \times_B U
}
$$
定义。它是 $U$ 上的群代数空间，其复合法则 $m$ 由复合法则
$c$ 诱导。
\end{lemma}

\begin{proof}
这是因为，在群胚范畴中，任一对象的自映射集构成一个群。
\end{proof}

\noindent
由于 $\Delta$ 是单态射，可知 $G = j^{-1}(\Delta_{U/B})$ 是 $R$ 的子层。
按这种方式看它时，结构态射 $G = j^{-1}(\Delta_{U/B}) \to U$ 由 $s$ 或
$t$ 诱导（两者相同），而 $m$ 由 $c$ 诱导。

\begin{definition}
\label{spaces-groupoids-definition-stabilizer-groupoid}
设 $B \to S$ 如第 \ref{spaces-groupoids-section-notation} 节所示。设
$(U, R, s, t, c)$ 为代数空间中的群胚（在 $B$ 上）。群代数空间
$j^{-1}(\Delta_{U/B}) \to U$ 称为emph{代数空间中的群胚
$(U, R, s, t, c)$ 的稳定子}。
\end{definition}

\noindent
文献中常用 $S$ 表示稳定子群代数空间（推测是因为 stabilizer 一词以 ``s''
开头）；我们不能这样做，因为已经用 $S$ 表示基概形。

\begin{lemma}
\label{spaces-groupoids-lemma-groupoid-action-stabilizer}
设 $B \to S$ 如第 \ref{spaces-groupoids-section-notation} 节所示。设
$(U, R, s, t, c)$ 为代数空间中的群胚（在 $B$ 上），并设 $G/U$ 为其稳定子。
用 $R_t/U$ 表示代数空间 $R$；它被看作 $U$ 上的代数空间，其结构态射为
$t : R \to U$。存在典范左作用
$$
a : G \times_U R_t \longrightarrow R_t
$$
它由复合法则 $c$ 诱导。
\end{lemma}

\begin{proof}
用 $T/B$ 上的点表示时，定义 $a(g, r) = c(g, r)$。
\end{proof}






\section{限制群胚}
\label{spaces-groupoids-section-restrict-groupoid}

\noindent
记号请参见《群胚》第 \ref{groupoids-section-restrict-groupoid} 节。

\begin{lemma}
\label{spaces-groupoids-lemma-restrict-groupoid}
设 $B \to S$ 如第 \ref{spaces-groupoids-section-notation} 节所示。设
$(U, R, s, t, c)$ 为代数空间中的群胚（在 $B$ 上）。设
$g : U' \to U$ 为代数空间的态射。考虑图表
$$
\xymatrix{
R' \ar[d] \ar[r] \ar@/_3pc/[dd]_{t'} \ar@/^1pc/[rr]^{s'}&
R \times_{s, U} U' \ar[r] \ar[d] &
U' \ar[d]^g \\
U' \times_{U, t} R \ar[d] \ar[r] &
R \ar[r]^s \ar[d]_t &
U \\
U' \ar[r]^g &
U
}
$$
其中所有方块都是纤维积方块。那么存在典范复合法则
$c' : R' \times_{s', U', t'} R' \to R'$，使得
$(U', R', s', t', c')$ 是代数空间中的群胚（在 $B$ 上），并且
$U' \to U$、$R' \to R$ 定义代数空间中的群胚的态射
$(U', R', s', t', c') \to (U, R, s, t, c)$（在 $B$ 上）。此外，对
任意概形 $T$（在 $B$ 上），群胚函子
$$
(U'(T), R'(T), s', t', c') \to (U(T), R(T), s, t, c)
$$
是 $(U(T), R(T), s, t, c)$ 经映射 $U'(T) \to U(T)$ 所作的限制
（见《群胚》第 \ref{groupoids-section-restrict-groupoid} 节）。
\end{lemma}

\begin{proof}
略。
\end{proof}

\begin{definition}
\label{spaces-groupoids-definition-restrict-groupoid}
设 $B \to S$ 如第 \ref{spaces-groupoids-section-notation} 节所示。设
$(U, R, s, t, c)$ 为代数空间中的群胚（在 $B$ 上）。设
$g : U' \to U$ 为 $B$ 上代数空间的态射。引理
\ref{spaces-groupoids-lemma-restrict-groupoid} 中构造的代数空间中的群胚的态射
$(U', R', s', t', c') \to (U, R, s, t, c)$ 称为
\emph{$(U, R, s, t, c)$ 到 $U'$ 的限制}。
% P09-ERR-0898
在这种情形下，有时使用记号 $R' = R|_{U'}$。
\end{definition}

\begin{lemma}
\label{spaces-groupoids-lemma-restrict-groupoid-relation}
定义 \ref{spaces-groupoids-definition-restrict-groupoid} 和
\ref{spaces-groupoids-definition-restrict-relation} 中所定义的限制群胚与限制（预）等价关系的概念，
通过引理 \ref{spaces-groupoids-lemma-groupoid-pre-equivalence} 和
\ref{spaces-groupoids-lemma-equivalence-groupoid} 的构造相符。
\end{lemma}

\begin{proof}
这里所说的是，引理 \ref{spaces-groupoids-lemma-restrict-groupoid} 中的 $R'$ 也等于
$$
R' = (U' \times_B U')\times_{U \times_B U} R
\longrightarrow
U' \times_B U'
$$
事实上，这也许是陈述该引理的一种更清楚的方式。
\end{proof}






\section{不变子空间}
\label{spaces-groupoids-section-invariant}

\noindent
本节简要讨论不变子空间的概念。

\begin{definition}
\label{spaces-groupoids-definition-invariant-open}
设 $B \to S$ 如第 \ref{spaces-groupoids-section-notation} 节所示。设
$(U, R, s, t, c)$ 为代数空间中的群胚（在基底 $B$ 上）。
\begin{enumerate}
\item 称开子空间 $W \subset U$ 是\emph{$R$-不变的}，如果
$t(s^{-1}(W)) \subset W$。
\item 若局部闭子空间 $Z \subset U$ 作为 $R$ 的局部闭子空间满足
$t^{-1}(Z) = s^{-1}(Z)$，则称它是\emph{$R$-不变的}。
\item 若代数空间的单态射 $T \to U$ 作为 $R$ 上的代数空间满足
$T \times_{U, t} R = R \times_{s, U} T$，则称它是\emph{$R$-不变的}。
\end{enumerate}
\end{definition}

\noindent
对于开子空间 $W \subset U$，$R$-不变性也等价于要求
$s^{-1}(W) = t^{-1}(W)$。如果 $W \subset U$ 是 $R$-不变的，那么
$R$ 到 $W$ 的限制就是 $R_W = s^{-1}(W) = t^{-1}(W)$。类似地，如果
$Z \subset U$ 是 $R$-不变局部闭子空间，那么 $R$ 到 $Z$ 的限制就是
$R_Z = s^{-1}(Z) = t^{-1}(Z)$。

\begin{lemma}
\label{spaces-groupoids-lemma-constructing-invariant-opens}
设 $B \to S$ 如第 \ref{spaces-groupoids-section-notation} 节所示。设
$(U, R, s, t, c)$ 为代数空间中的群胚（在 $B$ 上）。
\begin{enumerate}
\item 如果 $s$ 和 $t$ 都是开态射，那么对每个开集 $W \subset U$，开集
$s(t^{-1}(W))$ 是 $R$-不变的。
\item 如果 $s$ 和 $t$ 都开且拟紧，那么 $U$ 有一个开覆盖，其成员都是
$R$-不变拟紧开子空间。
\end{enumerate}
\end{lemma}

\begin{proof}
假设 $s$ 和 $t$ 都开，且 $W \subset U$ 开。由于 $s$ 开，可知
$W' = s(t^{-1}(W))$ 是 $U$ 的开子空间。
% P09-ERR-0899
从函子观点很容易看出这是 $R$-不变开子集（在 $U$ 中）；不过，我们认为直接使用一些
图表来论证颇有启发性。注意 $t^{-1}(W')$ 是态射
$$
A := t^{-1}(W) \times_{s|_{t^{-1}(W)}, U, t} R
\xrightarrow{\text{pr}_1} R
$$
的像，而 $s^{-1}(W')$ 是态射
$$
B := R \times_{s, U, s|_{t^{-1}(W)}} t^{-1}(W)
\xrightarrow{\text{pr}_0} R.
$$
上述箭头左侧的代数空间 $A$、$B$ 分别是
$R \times_{s, U, t} R$ 和 $R \times_{s, U, s} R$ 的开子空间。由引理
\ref{spaces-groupoids-lemma-diagram}，图表
$$
\xymatrix{
R \times_{s, U, t} R \ar[rd]_{\text{pr}_1} \ar[rr]_{(\text{pr}_1, c)} & &
R \times_{s, U, s} R \ar[ld]^{\text{pr}_0} \\
& R &
}
$$
交换，且水平箭头是同构。此外，显然
$(\text{pr}_1, c)(A) = B$。因此得到
$s^{-1}(W') = t^{-1}(W')$，并且 $W'$ 是 $R$-不变的。这证明了 (1)。

\medskip\noindent
现在假设 $s$、$t$ 都开且拟紧。那么，如果 $W \subset U$ 是拟紧开集，
$W' = s(t^{-1}(W))$ 也是拟紧开集，并且由上面的讨论可知它不变。让 $W$
遍历所有 \'etale 于 $U$ 的仿射概形的像，便得到 (2)。
\end{proof}






\section{商层}
\label{spaces-groupoids-section-quotient-sheaves}

\noindent
设 $S$ 为概形，$B$ 为 $S$ 上的代数空间。设
$j : R \to U \times_B U$ 为 $B$ 上的预关系。对每个概形 $S'$（在 $S$ 上），
可以取等价关系 $\sim_{S'}$，它由
$j(S') : R(S') \to U(S') \times U(S')$ 的像生成。于是得到预层
\begin{equation}
\label{spaces-groupoids-equation-quotient-presheaf}
\begin{matrix}
(\Sch/S)^{opp}_{fppf} &
\longrightarrow &
\textit{Sets}, \\
S' &
\longmapsto  &
U(S')/\sim_{S'}
\end{matrix}
\end{equation}
注意，因为 $j$ 是 $B$ 上代数空间的态射，且取值于 $U \times_B U$，
所以存在从预层（\ref{spaces-groupoids-equation-quotient-presheaf}）到 $B$ 的预层的典范变换。

\begin{definition}
\label{spaces-groupoids-definition-quotient-sheaf}
设 $B \to S$ 以及预关系 $j : R \to U \times_B U$ 如上。在此情形下，
\emph{商层 $U/R$}（与 $j$ 相伴）是预层
（\ref{spaces-groupoids-equation-quotient-presheaf}）在 $(\Sch/S)_{fppf}$ 上的层化。如果
$j : R \to U \times_B U$ 来自群代数空间 $G$（在 $B$ 上）对 $U$ 的作用，
如引理 \ref{spaces-groupoids-lemma-groupoid-from-action} 所述，那么把商层记为 $U/G$。
\end{definition}

\noindent
这恰好意味着图表
$$
\xymatrix{
R \ar@<1ex>[r] \ar@<-1ex>[r] &
U \ar[r] &
U/R
}
$$
是 $(\Sch/S)_{fppf}$ 上集合层范畴中的余等化子图表。仍然存在层的典范映射
$U/R \to B$，因为 $j$ 是 $B$ 上代数空间的态射，且取值于
$U \times_B U$。

\begin{remark}
\label{spaces-groupoids-remark-quotient-variant}
上述构造的一种变体是对函子
$$
\begin{matrix}
(\textit{Spaces}/B)^{opp}_{fppf} &
\longrightarrow &
\textit{Sets}, \\
X &
\longmapsto  &
U(X)/\sim_X
\end{matrix}
$$
作层化；这里 $\sim_X \subset U(X) \times U(X)$ 是由
$j : R(X) \to U(X) \times U(X)$ 的像生成的等价关系。当然，
$U(X) = \Mor_B(X, U)$ 且 $R(X) = \Mor_B(X, R)$。事实上，通过如下恒等，
所得结果相同：
% P09-ERR-0900
（在此插入《空间的拓扑》中的未来引用）。
\end{remark}

\begin{definition}
\label{spaces-groupoids-definition-representable-quotient}
% P09-ERR-0901
在定义 \ref{spaces-groupoids-definition-quotient-sheaf} 的情形下，称预关系 $j$ 有一个
\emph{可由代数空间表示的商}，如果层 $U/R$ 是代数空间。称预关系 $j$
有一个\emph{可表示商}，如果层 $U/R$ 可由概形表示。如果群胚
$(U, R, s, t, c)$ 是代数空间中的群胚（在 $B$ 上），并且取
商 $U/R$（其中 $j = (t, s)$）可表示（相应地，是代数空间），则称它有一个
\emph{可表示商}（相应地，\emph{可由代数空间表示的商}）。
\end{definition}

\noindent
如果商 $U/R$ 可由 $M$ 表示（表示对象或为概形，或为 $S$ 上的代数空间），
那么如上所见，它带有典范结构态射 $M \to B$。

\medskip\noindent
下一个引理刻画表示该商的 $M$。例如，当 $U \to M$ 平坦、有限表现且满，
并且 $R \cong U \times_M U$ 时，可以应用该引理。

\begin{lemma}
\label{spaces-groupoids-lemma-criterion-quotient-representable}
在定义 \ref{spaces-groupoids-definition-quotient-sheaf} 的情形下，假设存在代数空间
$M$（在 $S$ 上）和态射 $U \to M$，使得
\begin{enumerate}
\item 态射 $U \to M$ 等化 $s,t$，
\item 映射 $U \to M$ 是层的满射，并且
\item 诱导映射 $(t, s) : R \to U \times_M U$ 是层的满射。
\end{enumerate}
在这种情形下，$M$ 表示商层 $U/R$。
\end{lemma}

\begin{proof}
条件 (1) 说明 $U \to M$ 分解经过 $U/R$。条件 (2) 说明
$U/R \to M$ 作为层的映射是满射。条件 (3) 说明
$U/R \to M$ 作为层的映射是单射。因此引理成立。
\end{proof}

\noindent
如果不要求 $j$ 为预等价关系（例如，只要求它为预关系），下一个引理就是错的。

\begin{lemma}
\label{spaces-groupoids-lemma-quotient-pre-equivalence}
设 $S$ 为概形。设 $B$ 为 $S$ 上的代数空间。设
$j : R \to U \times_B U$ 为 $B$ 上的预等价关系。对概形 $S'$（在 $S$ 上）
和 $a, b \in U(S')$，下列条件等价：
\begin{enumerate}
\item $a$ 和 $b$ 映到 $(U/R)(S')$ 的同一个元素，并且
\item 存在 fppf 覆盖 $\{f_i : S_i \to S'\}$（覆盖 $S'$）和态射
$r_i : S_i \to R$，使得
$a \circ f_i = s \circ r_i$ 且 $b \circ f_i = t \circ r_i$。
\end{enumerate}
换言之，在此情形下，层的映射
$$
R \longrightarrow U \times_{U/R} U
$$
是满射。
\end{lemma}

\begin{proof}
略。提示：这个结论成立的原因是，在此情形下预层
（\ref{spaces-groupoids-equation-quotient-presheaf}）实际上由
$T \mapsto U(T)/j(R(T))$ 给出，因为
$j(R(T)) \subset U(T) \times U(T)$ 是等价关系；见定义
\ref{spaces-groupoids-definition-equivalence-relation}。
\end{proof}

\begin{lemma}
\label{spaces-groupoids-lemma-quotient-pre-equivalence-relation-restrict}
设 $S$ 为概形。设 $B$ 为 $S$ 上的代数空间。设
$j : R \to U \times_B U$ 为 $B$ 上的预等价关系，并设
$g : U' \to U$ 为 $B$ 上代数空间的态射。设
$j' : R' \to U' \times_B U'$ 为 $j$ 到 $U'$ 的限制。商层的映射
$$
U'/R' \longrightarrow U/R
$$
是单射。如果 $U' \to U$ 作为层的映射是满射，例如，如果
$\{g : U' \to U\}$ 是 fppf 覆盖（见《空间上的拓扑》定义
\ref{spaces-topologies-definition-fppf-covering}），那么
$U'/R' \to U/R$ 是层的同构。
\end{lemma}

\begin{proof}
假设截面 $\xi, \xi' \in (U'/R')(S')$ 映到 $U/R$ 的同一个截面。那么可以找到
fppf 覆盖 $\mathcal{S} = \{S_i \to S'\}$（覆盖 $S'$），使得
$\xi|_{S_i}, \xi'|_{S_i}$ 由 $a_i, a_i' \in U'(S_i)$ 给出。由引理
\ref{spaces-groupoids-lemma-quotient-pre-equivalence} 和位点公理，在作加细后可以假设
% P09-ERR-0902
$\mathcal{T}$ 存在态射 $r_i : S_i \to R$，使得
$g \circ a_i = s \circ r_i$、$g \circ a_i' = t \circ r_i$。由于根据构造
$R' = R \times_{U \times_S U} (U' \times_S U')$，可知
$(r_i, (a_i, a_i')) \in R'(S_i)$；这说明 $a_i$ 和 $a_i'$ 定义
$U'/R'$ 的同一个截面（在 $S_i$ 上）。由层条件，这推出 $\xi = \xi'$。

\medskip\noindent
如果 $U' \to U$ 是层的满射，那么 $U'/R' \to U/R$ 也是满射。最后，如果
$\{g : U' \to U\}$ 是 fppf 覆盖，那么层的映射 $U' \to U$ 是满射；见
《空间上的拓扑》引理 \ref{spaces-topologies-lemma-fppf-covering-surjective}。
\end{proof}

\begin{lemma}
\label{spaces-groupoids-lemma-quotient-groupoid-restrict}
设 $S$ 为概形。设 $B$ 为 $S$ 上的代数空间。设
$(U, R, s, t, c)$ 为代数空间中的群胚（在 $B$ 上）。
% P09-ERR-0904
设 $g : U' \to U$ 为 $B$ 上代数空间的态射。设
$(U', R', s', t', c')$ 为 $(U, R, s, t, c)$ 到 $U'$ 的限制。商层的映射
$$
U'/R' \longrightarrow U/R
$$
是单射。如果复合
$$
\xymatrix{
U' \times_{g, U, t} R \ar[r]_-{\text{pr}_1} \ar@/^3ex/[rr]^h
& R \ar[r]_s & U
}
$$
是 fppf 层的满射，那么该映射是双射。例如，当
$\{h : U' \times_{g, U, t} R \to U\}$ 是 $fppf$-覆盖，或
$U' \to U$ 是层的满射，或 $\{g : U' \to U\}$ 是 fppf 拓扑中的覆盖时，
这个条件成立。
\end{lemma}

\begin{proof}
单射性由引理 \ref{spaces-groupoids-lemma-groupoid-pre-equivalence} 和
\ref{spaces-groupoids-lemma-quotient-pre-equivalence-relation-restrict} 合并得到。为说明满射性
（层的满射映射的刻画见《位点》第 \ref{sites-section-sheaves-injective} 节），
按如下方式论证。假设 $T$ 为概形且 $\sigma \in U/R(T)$。存在覆盖
$\{T_i \to T\}$，使得 $\sigma|_{T_i}$ 是某个元素 $f_i \in U(T_i)$ 的像。
% P09-ERR-0905
因此，可以假设 $\sigma$ 是 $f \in U(T)$ 的像。由 $h$ 是层的满射这一假设，
可以找到 fppf 覆盖 $\{\varphi_i : T_i \to T\}$ 和态射
$f_i : T_i \to U' \times_{g, U, t} R$，使得
$f \circ \varphi_i = h \circ f_i$。记
$f'_i = \text{pr}_0 \circ f_i : T_i \to U'$。于是可见，
$f'_i \in U'(T_i)$ 映到 $g \circ f'_i \in U(T_i)$，并且
% P09-ERR-0906
$g \circ f'_i \sim_{T_i} h \circ f_i = f \circ \varphi_i$；记号如
（\ref{spaces-groupoids-equation-quotient-presheaf}）所示。给出该关系的 $R(T_i)$ 中元素是
$\text{pr}_1 \circ f_i$。这意味着，$\sigma$ 到 $T_i$ 的限制属于映射
% P09-ERR-0907
$U'/R'(T_i) \to U/R(T_i)$ 的像，正如所需。

\medskip\noindent
如果 $\{h\}$ 是 fppf 覆盖，那么它诱导层的满射；见《空间上的拓扑》引理
\ref{spaces-topologies-lemma-fppf-covering-surjective}。如果 $U' \to U$ 满，
那么 $h$ 也满，因为 $s$ 有一个截面
% P09-ERR-0908
（即群胚概形的中性元 $e$）。
\end{proof}






\section{商叠}
\label{spaces-groupoids-section-stacks}

\noindent
本节和随后的几节描述上面第 \ref{spaces-groupoids-section-quotient-sheaves} 节及《群胚》第
\ref{groupoids-section-quotient-sheaves} 节的一种推广。区别如下：我们将取商叠，
而不是商层。

\medskip\noindent
假设有概形 $S$、代数空间 $B$（在 $S$ 上），以及代数空间中的群胚
$(U, R, s, t, c)$（在 $B$ 上）。给定这些数据，考虑函子
\begin{equation}
\label{spaces-groupoids-equation-quotient-stack}
\begin{matrix}
(\Sch/S)_{fppf}^{opp} &
\longrightarrow &
\textit{Groupoids} \\
S' &
\longmapsto &
(U(S'), R(S'), s, t, c)
\end{matrix}
\end{equation}
由《范畴》例 \ref{categories-example-functor-groupoids}，这个“群胚预层”对应于
$(\Sch/S)_{fppf}$ 上的群胚纤维化范畴。在本章中把它记为
$$
[U/_{\!p}R] \to (\Sch/S)_{fppf}
$$
其中下标 ${}_p$ 用来与商叠相区别。

\begin{definition}
\label{spaces-groupoids-definition-quotient-stack}
商叠。设 $B \to S$ 如上。
\begin{enumerate}
\item 设 $(U, R, s, t, c)$ 为代数空间中的群胚（在 $B$ 上）。
\emph{商叠}
$$
p : [U/R] \longrightarrow (\Sch/S)_{fppf}
$$
属于 $(U, R, s, t, c)$；它是群胚纤维化范畴
$[U/_{\!p}R]$（在 $(\Sch/S)_{fppf}$ 上）的叠化（见《叠》引理
\ref{stacks-lemma-stackify-groupoids}），该范畴与
（\ref{spaces-groupoids-equation-quotient-stack}）相伴。
\item 设 $(G, m)$ 为 $B$ 上的群代数空间。设
% P09-ERR-0909
$a : G \times_B X \to X$ 为 $G$ 在 $B$ 上某个代数空间上的作用。
\emph{商叠}
$$
p : [X/G] \longrightarrow (\Sch/S)_{fppf}
$$
是与引理 \ref{spaces-groupoids-lemma-groupoid-from-action} 中代数空间中的群胚
$(X, G \times_B X, s, t, c)$（在 $B$ 上）相伴的商叠。
\end{enumerate}
\end{definition}

\noindent
因此 $[U/R]$ 和 $[X/G]$ 是 $(\Sch/S)_{fppf}$ 上的群胚叠。这些叠以后非常重要，
因而值得详细描述。回忆：给定代数空间 $X$（在 $S$ 上），我们用记号
$\mathcal{S}_X \to (\Sch/S)_{fppf}$ 表示与层 $X$ 相伴的集合叠；见《范畴》引理
\ref{categories-lemma-2-category-fibred-sets} 和《叠》引理
\ref{stacks-lemma-when-stack-in-sets}。

\begin{lemma}
\label{spaces-groupoids-lemma-quotient-stack-arrows}
假设 $B \to S$ 和 $(U, R, s, t, c)$ 如定义
\ref{spaces-groupoids-definition-quotient-stack} (1) 所示。存在典范 $1$-态射
$\pi : \mathcal{S}_U \to [U/R]$ 和 $[U/R] \to \mathcal{S}_B$，它们是
$(\Sch/S)_{fppf}$ 上群胚叠的态射。复合
$\mathcal{S}_U \to \mathcal{S}_B$ 是 $1$-态射，即与结构态射 $U \to B$ 相伴者。
\end{lemma}

\begin{proof}
在本证明中，用 $[U/_{\!p}R]$ 表示与群胚预层
（\ref{spaces-groupoids-equation-quotient-stack}）相伴的群胚纤维化范畴。由叠化的构造，存在
$1$-态射 $[U/_{\!p}R] \to [U/R]$。$1$-态射
$\mathcal{S}_U \to [U/R]$ 就是复合
$\mathcal{S}_U \to [U/_{\!p}R] \to [U/R]$；其中第一个箭头把概形
$S'/S$ 以及态射 $x : S' \to U$（在 $S$ 上）映到对象
$x \in U(S')$，它属于 $[U/_{\!p}R]$ 在 $S'$ 上的纤维范畴。

\medskip\noindent
为构造 $1$-态射 $[U/R] \to \mathcal{S}_B$，只需构造 $1$-态射
$[U/_{\!p}R] \to \mathcal{S}_B$；见《叠》引理
\ref{stacks-lemma-stackify-groupoids-universal-property}。在 $S'/S$ 上的对象处，
只使用映射
$$
U(S') \longrightarrow B(S')
$$
它来自结构态射 $U \to B$。显然，如果 $a \in R(S')$ 是一个“箭头”，其源为
$s(a) \in U(S')$、靶为 $t(a) \in U(S')$，那么因为 $s$ 和 $t$ 都是
\emph{$B$ 上}的态射，两者都映到同一个元素 $\overline{a}$（属于 $B(S')$）。
因此可将箭头 $a \in R(S')$ 映到 $\overline{a}$ 的恒等态射。（这是合适的，
因为纤维范畴 $(\mathcal{S}_B)_{S'}$ 只包含恒等态射。）我们略去验证这条规则与
这些分裂纤维化范畴上的拉回相容，因而如所需定义了 $1$-态射
$[U/_{\!p}R] \to \mathcal{S}_B$。

\medskip\noindent
略去最后一个断言的验证。
\end{proof}

\begin{lemma}
\label{spaces-groupoids-lemma-quotient-stack-2-arrow}
假设与记号如引理 \ref{spaces-groupoids-lemma-quotient-stack-arrows} 所示。存在典范 $2$-态射
$\alpha : \pi \circ s \to \pi \circ t$，它使图表
$$
\xymatrix{
\mathcal{S}_R \ar[r]_s \ar[d]_t & \mathcal{S}_U \ar[d]^\pi \\
\mathcal{S}_U \ar[r]^-\pi & [U/R]
}
$$
$2$-交换。
\end{lemma}

\begin{proof}
设 $S'$ 为 $S$ 上的概形。设 $r : S' \to R$ 为 $S$ 上的态射。那么
$r \in R(S')$ 是对象 $s \circ r, t \circ r \in U(S')$ 之间的同构。此外，
这个构造与拉回相容。这给出典范 $2$-态射
$\alpha_p : \pi_p \circ s \to \pi_p \circ t$，其中
$\pi_p : \mathcal{S}_U \to [U/_{\!p}R]$ 如引理
\ref{spaces-groupoids-lemma-quotient-stack-arrows} 的证明所示。因此，甚至图表
$$
\xymatrix{
\mathcal{S}_R \ar[r]_s \ar[d]^t & \mathcal{S}_U \ar[d]^{\pi_p} \\
\mathcal{S}_U \ar[r]^-{\pi_p} & [U/_{\!p}R]
}
$$
也是 $2$-交换的。因此更不必说，引理中的图表是 $2$-交换的。
\end{proof}

\begin{remark}
\label{spaces-groupoids-remark-fundamental-square}
在以后的各章中，我们将使用一种有歧义的记号：与其用
$\mathcal{S}_X$ 表示与 $X$ 相伴的集合叠，不如直接写成 $X$。使用这个记号，
引理 \ref{spaces-groupoids-lemma-quotient-stack-2-arrow} 的图表变成大家熟悉的图表
$$
\xymatrix{
R \ar[r]_s \ar[d]_t & U \ar[d]^\pi \\
U \ar[r]^-\pi & [U/R]
}
$$
在随后各节中，我们将证明这个图表有许多良好性质。特别地，将证明它是
$2$-纤维积（第 \ref{spaces-groupoids-section-quotient-stack-2-cartesian} 节），并且它几乎是
$2$-余等化子（对 $s$ 和 $t$ 而言；见第
\ref{spaces-groupoids-section-quotient-stacks-2-coequalize} 节）。
\end{remark}





\section{商叠的函子性}
\label{spaces-groupoids-section-functoriality-quotient-stacks}

\noindent
代数空间中的群胚的态射给出相伴的商叠态射。

\begin{lemma}
\label{spaces-groupoids-lemma-quotient-stack-functorial}
设 $S$ 为概形。设 $B$ 为 $S$ 上的代数空间。设
$f : (U, R, s, t, c) \to (U', R', s', t', c')$ 为代数空间中的群胚的态射
（在 $B$ 上）。那么 $f$ 诱导商叠的典范 $1$-态射
$$
[f] : [U/R] \longrightarrow [U'/R'].
$$
\end{lemma}

\begin{proof}
用 $[U/_{\!p}R]$ 和 $[U'/_{\!p}R']$ 表示与函子
（\ref{spaces-groupoids-equation-quotient-stack}）相伴的群胚纤维化范畴；它们取在基位点
$(\Sch/S)_{fppf}$ 上。显然，$f$ 定义 $1$-态射
$[U/_{\!p}R] \to [U'/_{\!p}R']$；把它与 $[U'/R']$ 的叠化函子复合，得到
$[U/_{\!p}R] \to [U'/R']$。
% P09-ERR-0910
再由叠化函子 $[U/_{\!p}R] \to [U/R]$ 的泛性质（见《叠》引理
\ref{stacks-lemma-stackify-groupoids-universal-property}），得到
$[U/R] \to [U'/R']$。
\end{proof}

\noindent
设 $B \to S$ 和 $f : (U, R, s, t, c) \to (U', R', s', t', c')$ 如引理
\ref{spaces-groupoids-lemma-quotient-stack-functorial} 所示。在此情形下，如下定义第三个代数空间中的
群胚（在 $B$ 上）；我们使用 $T$-值点的语言，其中 $T$ 是 $B$ 上变动的概形：
\begin{enumerate}
\item $U'' = U \times_{f, U', t'} R'$，所以一个 $T$-值点是偶
$(u, r')$，满足 $f(u) = t'(r')$；
\item $R'' = R \times_{f \circ s, U', t'} R'$，所以一个 $T$-值点是偶
$(r, r')$，满足 $f(s(r)) = t'(r')$；
\item $s'' : R'' \to U''$ 由 $s''(r, r') = (s(r), r')$ 给出；
\item $t'' : R'' \to U''$ 由 $t''(r, r') = (t(r), c'(f(r), r'))$ 给出；
\item $c'' : R'' \times_{s'', U'', t''} R'' \to R''$ 由
$c''((r_1, r'_1), (r_2, r'_2)) = (c(r_1, r_2), r'_2)$ 给出。
\end{enumerate}
$c''$ 的公式有意义，因为 $s''(r_1, r'_1) = t''(r_2, r'_2)$。显然
$c''$ 满足结合律。恒等元 $e''$ 由
% P09-ERR-0914
$e''(u, r) = (e(u), r)$ 给出。$(r, r')$ 的逆元由
$(i(r), c'(f(r), r'))$ 给出。因此，确实得到代数空间中的群胚（在 $B$ 上）。

\medskip\noindent
显然，映射 $U'' \to U$ 和 $R'' \to R$ 定义代数空间中的群胚的态射
$g : (U'', R'', s'', t'', c'') \to (U, R, s, t, c)$（在 $B$ 上）。此外，
映射 $U'' \to U'$、$(u, r') \mapsto s'(r')$ 以及
$R'' \to U'$、$(r, r') \mapsto s'(r')$ 表明，事实上
$(U'', R'', s'', t'', c'')$ 是代数空间中的群胚（在 $U'$ 上）。

\begin{lemma}
\label{spaces-groupoids-lemma-cartesian-square-of-morphism}
记号与假设如引理 \ref{spaces-groupoids-lemma-quotient-stack-functorial} 所示。设
$(U'', R'', s'', t'', c'')$ 为上面构造的代数空间中的群胚（在 $B$ 上）。
存在 $2$-交换方块
$$
\xymatrix{
[U''/R''] \ar[d] \ar[r]_{[g]} & [U/R] \ar[d]^{[f]} \\
\mathcal{S}_{U'} \ar[r] & [U'/R']
}
$$
它把 $[U''/R'']$ 识别为该 $2$-纤维积。
\end{lemma}

\begin{proof}
映射 $[f]$ 和 $[g]$ 来自应用引理
\ref{spaces-groupoids-lemma-quotient-stack-functorial}，另两个映射来自引理
\ref{spaces-groupoids-lemma-quotient-stack-arrows}（以及
$(U'', R'', s'', t'', c'')$ 取在 $U'$ 上这一事实）。为证明 $2$-纤维积性质，
只需对群胚纤维化范畴的图表
$$
\xymatrix{
[U''/_{\!p}R''] \ar[d] \ar[r]_{[g]} & [U/_{\!p}R] \ar[d]^{[f]} \\
\mathcal{S}_{U'} \ar[r] & [U'/_{\!p}R']
}
$$
证明引理；见《叠》引理
\ref{stacks-lemma-stackification-fibre-product-categories-fibred-in-groupoids}。
换言之，只需证明 $2$-纤维积的一个对象
% P09-ERR-0911
$\mathcal{S}_U \times_{[U'/_{\!p}R']} [U/_{\!p}R]$（在 $T$ 上）对应于
一个 $T$-值点（属于 $U''$），并对态射作类似证明。当然，这恰好是我们一开始构造
$U''$ 和 $R''$ 的方式。

\medskip\noindent
具体地，$\mathcal{S}_U \times_{[U'/_{\!p}R']} [U/_{\!p}R]$ 在 $T$ 上的
一个对象是三元组 $(u', u, r')$，其中 $u'$ 是 $T$-值点（属于 $U'$），$u$ 是
$T$-值点（属于 $U$），而 $r'$ 是从 $u'$ 到 $f(u)$ 的态射
% P09-ERR-0915
（在 $[U'/R']_T$ 中）；即
% P09-ERR-0912
$r'$ 是一个 $T$-值点（属于 $R$），满足 $s'(r') = u'$ 且
$t'(r') = f(u)$。显然，可以忘掉 $u'$ 而不丢失信息，并且可见这些对象与
% P09-ERR-0913
$T$-值点（属于 $R''$）一一对应。

\medskip\noindent
态射的情形类似：设 $(u'_1, u_1, r'_1)$ 和 $(u'_2, u_2, r'_2)$ 是该纤维积
在 $T$ 上的两个对象。那么，从 $(u'_2, u_2, r'_2)$ 到
$(u'_1, u_1, r'_1)$ 的态射由 $(1, r)$ 给出，其中
$1 : u'_1 \to u'_2$ 仅仅意味着 $u'_1 = u'_2$（这是因为
$\mathcal{S}_U$ 纤维化于集合），而 $r$ 是 $T$-值点（属于 $R$），满足
$s(r) = u_2$、$t(r) = u_1$，并且还有
$c'(f(r), r'_2) = r'_1$。因此箭头
$$
(1, r) : (u'_2, u_2, r'_2) \to (u'_1, u_1, r'_1)
$$
完全由偶 $(r, r'_2)$ 决定。于是箭头函子由 $R''$ 表示，而且态射
$s''$、$t''$、$c''$ 显然分别对应于 $2$-纤维积
$\mathcal{S}_U \times_{[U'/_{\!p}R']} [U/_{\!p}R]$ 中的源、靶与复合。
\end{proof}






\section{商叠的 2-笛卡儿方块}
\label{spaces-groupoids-section-quotient-stack-2-cartesian}

\noindent
本节计算商叠的 $\mathit{Isom}$-层，并由此推出商叠的定义图表是
$2$-纤维积。

\begin{lemma}
\label{spaces-groupoids-lemma-quotient-stack-morphisms}
假设 $B \to S$、$(U, R, s, t, c)$ 和
$\pi : \mathcal{S}_U \to [U/R]$ 如引理
\ref{spaces-groupoids-lemma-quotient-stack-arrows} 所示。设 $S'$ 为 $S$ 上的概形。设
$x, y \in \Ob([U/R]_{S'})$ 为商叠在 $S'$ 上的对象。如果
$x = \pi(x')$ 且 $y = \pi(y')$，其中 $x', y' : S' \to U$ 是某些态射，那么
$$
\mathit{Isom}(x, y) = S' \times_{(y', x'), U \times_S U} R
$$
作为 $S'$ 上的层。
\end{lemma}

\begin{proof}
设 $[U/_{\!p}R]$ 为与群胚预层（\ref{spaces-groupoids-equation-quotient-stack}）相伴的
群胚纤维化范畴，如引理 \ref{spaces-groupoids-lemma-quotient-stack-arrows} 的证明所述。由构造，
层 $\mathit{Isom}(x, y)$ 是与预层 $\mathit{Isom}(x', y')$ 相伴的层。另一方面，
由 $[U/_{\!p}R]$ 中态射的定义，有
$$
\mathit{Isom}(x', y') = S' \times_{(y', x'), U \times_S U} R
$$
而右侧是代数空间，因而是层。
\end{proof}

\begin{lemma}
\label{spaces-groupoids-lemma-quotient-stack-2-cartesian}
假设 $B \to S$、$(U, R, s, t, c)$ 和
$\pi : \mathcal{S}_U \to [U/R]$ 如引理
\ref{spaces-groupoids-lemma-quotient-stack-arrows} 所示。引理
\ref{spaces-groupoids-lemma-quotient-stack-2-arrow} 的 $2$-交换方块
$$
\xymatrix{
\mathcal{S}_R \ar[r]_s \ar[d]_t & \mathcal{S}_U \ar[d]^\pi \\
\mathcal{S}_U \ar[r]^-\pi & [U/R]
}
$$
是一个 $2$-纤维积（在 $(\Sch/S)_{fppf}$ 的群胚叠中）。
\end{lemma}

\begin{proof}
由《叠》引理 \ref{stacks-lemma-2-product-stacks-in-groupoids}，该引理的断言
有意义。它还说明，我们必须证明函子
$$
\mathcal{S}_R \longrightarrow \mathcal{S}_U \times_{[U/R]} \mathcal{S}_U
$$
是等价；该函子把 $r : T \to R$ 映到
$(T, t(r), s(r), \alpha(r))$，其中右侧是《范畴》引理
\ref{categories-lemma-2-product-categories-over-C} 中所述的 $2$-纤维积。
把定义展开以后，这恰好是引理
\ref{spaces-groupoids-lemma-quotient-stack-morphisms} 的内容。（另一种证明：
% P09-ERR-0916
在此情形下展开引理 \ref{spaces-groupoids-lemma-cartesian-square-of-morphism} 的含义，
也会得到该结论。）
\end{proof}

\begin{lemma}
\label{spaces-groupoids-lemma-quotient-stack-isom}
假设 $B \to S$ 和 $(U, R, s, t, c)$ 如定义
\ref{spaces-groupoids-definition-quotient-stack} (1) 所示。对任意概形 $T$（在 $S$ 上）以及
对象 $x,y$（属于 $[U/R]$，在 $T$ 上），层
$\mathit{Isom}(x, y)$（在 $(\Sch/T)_{fppf}$ 上）具有如下性质：存在 fppf 覆盖
$\{T_i \to T\}_{i \in I}$，使得
$\mathit{Isom}(x, y)|_{(\Sch/T_i)_{fppf}}$ 可由代数空间表示。
\end{lemma}

\begin{proof}
这直接由引理 \ref{spaces-groupoids-lemma-quotient-stack-morphisms} 以及如下事实得到：根据商叠的
构造，$x$ 和 $y$ 在 fppf 拓扑下局部地都来自 $\mathcal{S}_U$ 的对象。
\end{proof}






\section{商叠的 2-余等化子性质}
\label{spaces-groupoids-section-quotient-stacks-2-coequalize}

\noindent
在群胚上有复合，它为上一个引理中的典范 $2$-态射导出上闭链条件。为了给出精确
表述，我们使用《范畴》第 \ref{categories-section-formal-cat-cat} 节和第
\ref{categories-section-2-categories} 节中引入的记号。

\begin{lemma}
\label{spaces-groupoids-lemma-quotient-stack-cocycle}
假设与记号如引理 \ref{spaces-groupoids-lemma-quotient-stack-arrows} 和
\ref{spaces-groupoids-lemma-quotient-stack-2-arrow} 所示。下图中两个 2-态射的竖直复合
$$
\xymatrix@C=15pc{
\mathcal{S}_{R \times_{s, U, t} R}
\ruppertwocell^{\pi \circ s \circ \text{pr}_1 = \pi \circ s \circ c}{\ \ \ \ \ \ \alpha \star \text{id}_{\text{pr}_1}}
\ar[r]_(.3){\pi \circ t \circ \text{pr}_1 = \pi \circ s \circ \text{pr}_0}
\rlowertwocell_{\pi \circ t \circ \text{pr}_0 = \pi \circ t \circ c}{\ \ \ \ \ \ \alpha \star \text{id}_{\text{pr}_0}}
&
[U/R]
}
$$
是 $2$-态射 $\alpha \star \text{id}_c$。用公式表示，
$\alpha \star \text{id}_c =
(\alpha \star \text{id}_{\text{pr}_0})
\circ
(\alpha \star \text{id}_{\text{pr}_1})
$。
\end{lemma}

\begin{proof}
作两点说明：
\begin{enumerate}
\item 公式
$\alpha \star \text{id}_c = (\alpha \star \text{id}_{\text{pr}_0}) \circ
(\alpha \star \text{id}_{\text{pr}_1})$ 仅在实现下列\emph{等式}后才有意义：
$\pi \circ s \circ \text{pr}_1 = \pi \circ s \circ c$、
$\pi \circ t \circ \text{pr}_1 = \pi \circ s \circ \text{pr}_0$ 和
$\pi \circ t \circ \text{pr}_0 = \pi \circ t \circ c$。确切地说，第二个等式
保证竖直复合 $\circ$ 有意义，另两个等式则保证公式两侧都是有相同源与靶的
$2$-态射。
\item 引理成立的原因是，与群胚预层（\ref{spaces-groupoids-equation-quotient-stack}）相伴的
群胚纤维化范畴 $[U/_{\!p}R]$ 中的复合来自复合法则
$c : R \times_{s, U, t} R \to R$。
\end{enumerate}
略去引理的证明。
\end{proof}

\noindent
注意，在引理的情形下，实际上已有等式
$s \circ \text{pr}_1 = s \circ c$、
$t \circ \text{pr}_1 = s \circ \text{pr}_0$ 和
$t \circ \text{pr}_0 = t \circ c$；这些等式是在与 $\pi$ 复合之前成立的。因此，下一个引理中的公式与上一个引理中的
公式以完全相同的方式有意义。

\begin{lemma}
\label{spaces-groupoids-lemma-quotient-stack-2-coequalizer}
假设与记号如引理 \ref{spaces-groupoids-lemma-quotient-stack-arrows} 和
\ref{spaces-groupoids-lemma-quotient-stack-2-arrow} 所示。引理
\ref{spaces-groupoids-lemma-quotient-stack-2-arrow} 的 $2$-交换图表在如下意义下是
$2$-余等化子：给定
\begin{enumerate}
\item 群胚叠 $\mathcal{X}$（在 $(\Sch/S)_{fppf}$ 上）；
\item $1$-态射 $f : \mathcal{S}_U \to \mathcal{X}$；以及
\item $2$-箭头 $\beta : f \circ s \to f \circ t$，
\end{enumerate}
使得
$$
\beta \star \text{id}_c
=
(\beta \star \text{id}_{\text{pr}_0})
\circ
(\beta \star \text{id}_{\text{pr}_1})
$$
那么存在 $1$-态射 $[U/R] \to \mathcal{X}$，它使图表
$$
\xymatrix{
\mathcal{S}_R \ar[r]_s \ar[d]^t & \mathcal{S}_U \ar[d] \ar[ddr]^f \\
\mathcal{S}_U \ar[r] \ar[rrd]_f & [U/R] \ar[rd] \\
& & \mathcal{X}
}
$$
$2$-交换。
\end{lemma}

\begin{proof}
假设给定引理中的 $\mathcal{X}$、$f$ 和 $\beta$。由《叠》引理
\ref{stacks-lemma-stackify-groupoids-universal-property}，只需构造 $1$-态射
$g : [U/_{\!p}R] \to \mathcal{X}$。首先注意，$1$-态射
$\mathcal{S}_U \to [U/_{\!p}R]$ 在对象上是双射。因此在对象上可以置
$g(x) = f(x)$，其中
$x \in \Ob(\mathcal{S}_U) = \Ob([U/_{\!p}R])$。态射
$\varphi : x \to y$（属于 $[U/_{\!p}R]$）来自交换图表
$$
\xymatrix{
S_2 \ar[dd]_h \ar[r]_x \ar[dr]_\varphi & U \\
& R \ar[u]_s \ar[d]^t \\
S_1 \ar[r]^y & U.
}
$$
于是可以把 $g(\varphi)$ 置为复合
$$
\xymatrix{
f(x) \ar@{=}[r] \ar[rrrrrd] &
f(s \circ \varphi) \ar@{=}[r] &
(f \circ s)(\varphi) \ar[r]^\beta &
(f \circ t)(\varphi) \ar@{=}[r] &
f(t \circ \varphi) \ar@{=}[r] &
f(y \circ h) \ar[d] \\
& & & & & f(y).
}
$$
竖直箭头是把函子 $f$ 应用于典范态射
$y \circ h \to y$（在 $\mathcal{S}_U$ 中）所得的结果（即提升 $h$ 且以
$y$ 为靶的强笛卡儿态射）。
% P09-ERR-0917
我们验证如此定义的 $f$ 与复合相容，至少在纤维范畴上如此。设 $S'$ 为 $S$ 上的
概形，并设 $a : S' \to R \times_{s, U, t} R$ 为态射。在此情形下，置
$x = s \circ \text{pr}_1 \circ a = s \circ c \circ a$、
$y = t \circ \text{pr}_1 \circ a = s \circ \text{pr}_0 \circ a$，以及
% P09-ERR-0918
$z = t \circ \text{pr}_0 \circ a = t \circ \text{pr}_0 \circ c$，从而得到交换图表
$$
\xymatrix{
x \ar[rr]_{c \circ a} \ar[rd]_{\text{pr}_1 \circ a} & & z \\
& y \ar[ru]_{\text{pr}_0 \circ a}
}
$$
它位于纤维范畴 $[U/_{\!p}R]_{S'}$ 中。
此外，该纤维范畴中的每个交换三角形都具有这种形式。于是由上面的定义可知，
$f$ 把它映到交换图表，当且仅当图表
$$
\xymatrix{
& (f \circ s)(c \circ a) \ar[r]_-{\beta} &
(f \circ t)(c \circ a) \ar@{=}[rd] & \\
(f \circ s)(\text{pr}_1 \circ a) \ar[rd]^\beta \ar@{=}[ru] & & &
(f \circ t)(\text{pr}_0 \circ a) \\
& (f \circ t)(\text{pr}_1 \circ a) \ar@{=}[r] &
(f \circ s)(\text{pr}_0 \circ a) \ar[ru]^\beta
}
$$
交换；这恰好是引理中公式所表达的条件。略去验证 $f$ 把恒等态射映到恒等态射，
并且与任意态射的复合相容。
\end{proof}

\section{商叠的显式描述}
\label{spaces-groupoids-section-explicit-quotient-stacks}

\noindent
为陈述结果，我们需要引入一些记号。假设 $B \to S$ 以及
$(U, R, s, t, c)$ 如定义 \ref{spaces-groupoids-definition-quotient-stack} (1) 中所述。
设 $T$ 为 $S$ 上的概形。设
$\mathcal{T} = \{T_i \to T\}_{i \in I}$ 为一个 fppf 覆盖。一个
{\it $[U/R]$-下降资料}（相对于 $\mathcal{T}$）由系统
$(u_i, r_{ij})$ 给出，其中
\begin{enumerate}
\item 对每个 $i$，有一个态射 $u_i : T_i \to U$，并且
\item 对每个 $i, j$，有一个态射 $r_{ij} : T_i \times_T T_j \to R$，
\end{enumerate}
使得
\begin{enumerate}
\item[(a)] 作为态射 $T_i \times_T T_j \to U$，有
$$
s \circ r_{ij} = u_i \circ \text{pr}_0
\quad\text{和}\quad
t \circ r_{ij} = u_j \circ \text{pr}_1,
$$
\item[(b)] 作为态射 $T_i \times_T T_j \times_T T_k \to R$，有
$$
c \circ (r_{jk} \circ \text{pr}_{12}, r_{ij} \circ \text{pr}_{01})
= r_{ik} \circ \text{pr}_{02}.
$$
\end{enumerate}
{\it 态射 $(u_i, r_{ij}) \to (u'_i, r'_{ij})$} 是两个
$[U/R]$-下降资料在同一个覆盖 $\mathcal{T}$ 上的态射，指一个族
$(r_i : T_i \to R)$，使得
\begin{enumerate}
\item[] $(\alpha)$ 作为态射 $T_i \to U$，有
$$
u_i = s \circ r_i
\quad\text{和}\quad
u'_i = t \circ r_i
$$
\item[] $(\beta)$ 作为态射 $T_i \times_T T_j \to R$，有
$$
c \circ (r'_{ij}, r_i \circ \text{pr}_0)
=
c \circ (r_j \circ \text{pr}_1, r_{ij}).
$$
\end{enumerate}
相对于一个固定覆盖的下降资料态射上有自然的复合法则，由此得到下降资料范畴。
这个范畴是一个群胚。最后，若
$\mathcal{T}' = \{T'_j \to T\}_{j \in J}$
是加细 $\mathcal{T}$ 的第二个 fppf 覆盖，则下降资料存在拉回的概念。在这种情形下，
它尤其容易显式描述。确切地说，若 $\alpha : J \to I$ 且
% P09-ERR-0919
$\varphi_j : T'_j \to T_{\alpha(i)}$ 是覆盖的态射，则下降资料
$(u_i, r_{ii'})$ 的拉回就是
$$
(u_{\alpha(i)} \circ \varphi_j,
r_{\alpha(j)\alpha(j')} \circ \varphi_j \times \varphi_{j'}).
$$
% P09-ERR-0920
以这种方式定义的拉回给出一个从 $\mathcal{T}$ 上的下降资料范畴到
$\mathcal{T}'$ 上的下降资料范畴的函子。

\begin{lemma}
\label{spaces-groupoids-lemma-quotient-stack-objects}
假设 $B \to S$ 以及 $(U, R, s, t, c)$ 如
定义 \ref{spaces-groupoids-definition-quotient-stack} (1) 中所述。
设 $\pi : \mathcal{S}_U \to [U/R]$ 如
引理 \ref{spaces-groupoids-lemma-quotient-stack-arrows} 中所述。
设 $T$ 为 $S$ 上的概形。
\begin{enumerate}
\item 对每个对象 $x$，若它属于纤维范畴 $[U/R]_T$，则存在一个 fppf 覆盖
$\{f_i : T_i \to T\}_{i \in I}$，使得
$f_i^*x \cong \pi(u_i)$，其中 $u_i \in U(T_i)$；
\item 下列同构的复合
$$
\pi(u_i \circ \text{pr}_0)
=
\text{pr}_0^*\pi(u_i)
\cong
\text{pr}_0^*f_i^*x
\cong
\text{pr}_1^*f_j^*x
\cong
\text{pr}_1^*\pi(u_j)
=
\pi(u_j \circ \text{pr}_1)
$$
具有 $\pi(r_{ij})$ 的形式，其中
$r_{ij} : T_i \times_T T_j \to R$ 是某些态射；
\item 系统 $(u_i, r_{ij})$ 构成如上定义的一个 $[U/R]$-下降资料；
\item 任意 $[U/R]$-下降资料 $(u_i, r_{ij})$ 都以这种方式产生；
\item 若 $x$ 如上对应于 $(u_i, r_{ij})$，而
$y \in \Ob([U/R]_T)$ 对应于 $(u'_i, r'_{ij})$，则有典范双射
$$
\Mor_{[U/R]_T}(x, y)
\longleftrightarrow
\left\{
\begin{matrix}
\text{态射 }(u_i, r_{ij}) \to (u'_i, r'_{ij})\\
\text{作为 }[U/R]\text{-下降资料}
\end{matrix}
\right\}
$$
\item 这个对应与 fppf 覆盖的加细相容。
\end{enumerate}
\end{lemma}

\begin{proof}
% P09-ERR-0921
断言 (1) 是叠化构造的一部分。
第 (2) 部分由引理 \ref{spaces-groupoids-lemma-quotient-stack-morphisms} 得出。
略去对 (3) 的验证。第 (4) 部分是叠中所有下降资料均有效这一事实的改述。
略去对 (5) 和 (6) 的验证。
\end{proof}

\section{限制与商叠}
\label{spaces-groupoids-section-quotient-stack-restrict}

\noindent
本节研究取限制时商叠会发生什么。

\begin{lemma}
\label{spaces-groupoids-lemma-quotient-stack-restrict}
记号与假设如引理 \ref{spaces-groupoids-lemma-quotient-stack-functorial} 中所述。
商叠的态射
$$
[f] : [U/R] \longrightarrow [U'/R']
$$
全忠实，当且仅当 $R$ 是 $R'$ 通过态射 $f : U \to U'$ 得到的限制。
\end{lemma}

\begin{proof}
设 $x, y$ 是 $[U/R]$ 在概形 $T/S$ 上的对象。设 $x', y'$ 是 $x, y$
在范畴 $[U'/R']_T$ 中的像。函子 $[f]$ 全忠实，当且仅当层的态射
$$
\mathit{Isom}(x, y) \longrightarrow \mathit{Isom}(x', y')
$$
对每个 $T, x, y$ 都是同构。可以在 $T$ 上局部地（就 fppf 拓扑而言）检验这一点。
因此，由引理 \ref{spaces-groupoids-lemma-quotient-stack-objects}，可以假设 $x, y$ 来自
$a, b \in U(T)$。在这种情形下，可见 $x', y'$ 对应于
$f \circ a, f \circ b$。由引理 \ref{spaces-groupoids-lemma-quotient-stack-morphisms}，
在这种情形下，上面显示的层态射成为
$$
T \times_{(a, b), U \times_B U} R
\longrightarrow
T \times_{f \circ a, f \circ b, U' \times_B U'} R'.
$$
若 $R$ 是该限制，这就是同构，因为在这种情形下
$R = (U \times_B U) \times_{U' \times_B U'} R'$；参见
引理 \ref{spaces-groupoids-lemma-restrict-groupoid-relation} 及其证明。反过来，若最后显示的态射
对所有 $T, a, b$ 都是同构，则可推出
$R = (U \times_B U) \times_{U' \times_B U'} R'$，即
$R$ 是 $R'$ 的限制。
\end{proof}

\begin{lemma}
\label{spaces-groupoids-lemma-quotient-stack-restrict-equivalence}
记号与假设如引理 \ref{spaces-groupoids-lemma-quotient-stack-functorial} 中所述。
商叠的态射
$$
[f] : [U/R] \longrightarrow [U'/R']
$$
是等价，当且仅当
\begin{enumerate}
\item $(U, R, s, t, c)$ 是 $(U', R', s', t', c')$ 通过
$f : U \to U'$ 得到的限制，并且
\item 映射
$$
\xymatrix{
U \times_{f, U', t'} R' \ar[r]_-{\text{pr}_1} \ar@/^3ex/[rr]^h
& R' \ar[r]_{s'} & U'
}
$$
是层的满射。
\end{enumerate}
例如，若 $\{h : U \times_{f, U', t'} R' \to U'\}$ 是 fppf 覆盖，
或 $f : U \to U'$ 是层的满射，或 $\{f : U \to U'\}$ 是 fppf 覆盖，
则第 (2) 部分成立。
\end{lemma}

\begin{proof}
% P09-ERR-0922
由引理 \ref{spaces-groupoids-lemma-quotient-stack-restrict}，我们已经知道第 (1) 部分等价于
全忠实性。因此可以假设 (1) 成立且 $[f]$ 全忠实。我们的目标是在这些假设下证明：
$[f]$ 是等价，当且仅当 (2) 成立。可以使用刻画等价的
Stacks，引理 \ref{stacks-lemma-characterize-essentially-surjective-when-ff}。

\medskip\noindent
假设 (2)。我们将使用
Stacks，引理 \ref{stacks-lemma-characterize-essentially-surjective-when-ff}
来证明 $[f]$ 是等价。设 $T$ 是概形且
$x' \in \Ob([U'/R']_T)$。存在一个覆盖 $\{g_i : T_i \to T\}$，
使得 $g_i^*x'$ 是某个元素 $a'_i \in U'(T_i)$ 的像；参见
引理 \ref{spaces-groupoids-lemma-quotient-stack-objects}。因此可以假设 $x'$ 是
$a' \in U'(T)$ 的像。由 $h$ 是层的满射这一假设，可以找到一个 fppf 覆盖
$\{\varphi_i : T_i \to T\}$ 和态射
% P09-ERR-0923
$b_i : T_i \to U \times_{g, U', t'} R'$，使得
$a' \circ \varphi_i = h \circ b_i$。记
$a_i = \text{pr}_0 \circ b_i : T_i \to U$。于是可见
$a_i \in U(T_i)$ 映到 $f \circ a_i \in U'(T_i)$，并且
$f \circ a_i \cong_{T_i} h \circ b_i = a' \circ \varphi_i$，
其中 $\cong_{T_i}$ 表示纤维范畴 $[U'/R']_{T_i}$ 中的同构。确切地说，
给出该同构的 $R'(T_i)$ 中元素是 $\text{pr}_1 \circ b_i$。
% P09-ERR-0924
这意味着 $x$ 到 $T_i$ 的限制位于函子
$[U/R]_{T_i} \to [U'/R']_{T_i}$ 的本质像中，正如所需。

\medskip\noindent
假设 $[f]$ 是等价。令 $\xi' \in [U'/R']_{U'}$ 表示对应于 $U'$
的恒等态射的对象。应用
Stacks，引理 \ref{stacks-lemma-characterize-essentially-surjective-when-ff}，
可见存在一个 fppf 覆盖
$\mathcal{U}' = \{g'_i : U'_i \to U'\}$，使得
$(g'_i)^*\xi' \cong [f](\xi_i)$，其中某个 $\xi_i$ 位于
$[U/R]_{U'_i}$ 中。加细覆盖 $\mathcal{U}'$ 之后（使用
引理 \ref{spaces-groupoids-lemma-quotient-stack-objects}），可以假设 $\xi_i$ 来自态射
$a_i : U'_i \to U$。事实
$[f](\xi_i) \cong (g'_i)^*\xi'$ 意味着，在必要时再次加细覆盖
$\mathcal{U}'$ 后，存在态射 $r'_i : U'_i \to R'$，满足
$t' \circ r'_i = f \circ a_i$ 且
$s' \circ r'_i = \text{id}_{U'} \circ g'_i$。图示如下：
$$
\xymatrix{
U \ar[d]^f & & U'_i \ar[ll]^{a_i} \ar[ld]_{r'_i} \ar[d]^{g'_i} \\
U' & R' \ar[l]_{t'} \ar[r]^{s'} & U'
}
$$
因此 $(a_i, r'_i) : U'_i \to U \times_{g, U', t'} R'$ 是态射，
满足 $h \circ (a_i, r'_i) = g'_i$；由此可知
$\{h : U \times_{g, U', t'} R' \to U'\}$ 可由 fppf 覆盖
$\mathcal{U}'$ 加细，这意味着 $h$ 诱导层的满射；参见
代数空间上的拓扑，引理
\ref{spaces-topologies-lemma-fppf-covering-surjective}。

\medskip\noindent
若 $\{h\}$ 是 fppf 覆盖，则它诱导层的满射；参见
代数空间上的拓扑，引理
\ref{spaces-topologies-lemma-fppf-covering-surjective}。
% P09-ERR-0925
若 $U' \to U$ 是满射，则 $h$ 也是满射，因为 $s$ 有一个截面
（即代数空间群胚的中性元 $e$）。
\end{proof}

\begin{lemma}
\label{spaces-groupoids-lemma-criterion-fibre-product}
记号与假设如引理 \ref{spaces-groupoids-lemma-quotient-stack-functorial} 中所述。
假设
$$
\xymatrix{
R \ar[d]_s \ar[r]_f & R' \ar[d]^{s'} \\
U \ar[r]^f & U'
}
$$
是笛卡儿的。则
$$
\xymatrix{
\mathcal{S}_U \ar[d] \ar[r] & [U/R] \ar[d]^{[f]} \\
\mathcal{S}_{U'} \ar[r] & [U'/R']
}
$$
是一个 $2$-纤维积方块。
\end{lemma}

\begin{proof}
把逆同构 $i : R \to R$ 和 $i' : R' \to R'$ 应用于引理断言中的
（第一个）笛卡儿图表，可见
$$
\xymatrix{
R \ar[d]_t \ar[r]_f & R' \ar[d]^{t'} \\
U \ar[r]^f & U'
}
$$
也是笛卡儿的。由引理 \ref{spaces-groupoids-lemma-cartesian-square-of-morphism}，有一个
$2$-纤维积方块
$$
\xymatrix{
[U''/R''] \ar[d] \ar[r] & [U/R] \ar[d] \\
\mathcal{S}_{U'} \ar[r] & [U'/R']
}
$$
其中 $U'' = U \times_{f, U', t'} R'$ 且
$R'' = R \times_{f \circ s, U', t'} R'$。由上可见
$(t, f) : R \to U''$ 是同构，并且
% P09-ERR-0926
$$
R'' =
R \times_{f \circ s, U', t'} R' =
R \times_{s, U} U \times_{f, U', t'} R' =
R \times_{s, U, t} \times R.
$$
显式地说，同构 $R \times_{s, U, t} R \to R''$ 由法则
$(r_0, r_1) \mapsto (r_0, f(r_1))$ 给出。此外，$s'', t'', c''$
转化为映射
$$
R \times_{s, U, t} R \to R,
\quad
s''(r_0, r_1) = r_1, \quad t''(r_0, r_1) = c(r_0, r_1)
$$
以及
$$
\begin{matrix}
c'' : &
(R \times_{s, U, t} R) \times_{s'', R, t''} (R \times_{s, U, t} R)
&
\longrightarrow
&
R \times_{s, U, t} R, \\
&
((r_0, r_1), (r_2, r_3)) &
\longmapsto & (c(r_0, r_2), r_3).
\end{matrix}
$$
预复合以同构
$$
R \times_{s, U, s} R \longrightarrow R \times_{s, U, t} R,
\quad
(r_0, r_1) \longmapsto (c(r_0, i(r_1)), r_1)
$$
后，可见 $t''$ 和 $s''$ 转化为 $\text{pr}_0$ 和 $\text{pr}_1$，
而 $c''$ 转化为
$\text{pr}_{02} :
R \times_{s, U, s} R \times_{s, U, s} R \to R \times_{s, U, s} R$。
因此可见有同构
$[U''/R''] \cong [R/R \times_{s, U, s} R]$，其中，作为代数空间中的群胚，
$(R, R \times_{s, U, s} R, s'', t'', c'')$ 是平凡群胚
$(U, U, \text{id}, \text{id}, \text{id})$ 通过 $s : R \to U$ 得到的限制。
由于 $s : R \to U$ 是 fppf 层的满射（因为它有右逆），态射
$$
[U''/R''] \cong [R/R \times_{s, U, s} R]
\longrightarrow
[U/U] = \mathcal{S}_U
$$
由引理 \ref{spaces-groupoids-lemma-quotient-stack-restrict-equivalence} 是等价。这就证明了引理。
\end{proof}

\section{惯性与商叠}
\label{spaces-groupoids-section-inertia}

\noindent
群胚叠的（相对）惯性叠定义于
Stacks，第 \ref{stacks-section-the-inertia-stack} 节。
% P09-ERR-0927
在纤维化范畴的设定中，实际构造及其若干性质见
范畴，第 \ref{categories-section-inertia} 节。

\begin{lemma}
\label{spaces-groupoids-lemma-presentation-inertia}
假设 $B \to S$ 以及 $(U, R, s, t, c)$ 如
定义 \ref{spaces-groupoids-definition-quotient-stack} (1) 中所述。设 $G/U$ 为群胚
$(U, R, s, t, c, e, i)$ 的稳定子群代数空间；参见
定义 \ref{spaces-groupoids-definition-stabilizer-groupoid}。置 $R' = R \times_{s, U} G$，并置
\begin{enumerate}
\item $s' : R' \to G$，$(r, g) \mapsto g$；
\item $t' : R' \to G$，$(r, g) \mapsto c(r, c(g, i(r)))$；
% P09-ERR-0928
\item $c' : R' \times_{s', G, t'} R' \to R'$，
$((r_1, g_1), (r_2, g_2) \mapsto (c(r_1, r_2), g_1)$。
\end{enumerate}
则 $(G, R', s', t', c')$ 是 $B$ 上代数空间中的群胚，并且
$$
\mathcal{I}_{[U/R]} = [G/ R'].
$$
即，相伴商叠是 $[U/R]$ 的惯性叠。
\end{lemma}

\begin{proof}
由 Stacks，引理 \ref{stacks-lemma-stackification-inertia}，只需证明
$\mathcal{I}_{[U/_{\!p}R]} = [G/_{\!p} R']$。设 $T$ 为 $S$ 上的概形。
回顾，$[U/_{\!p}R]$ 在 $T$ 上的惯性纤维化范畴的一个对象由一对
$(x, g)$ 给出，其中 $x$ 是 $[U/_{\!\!p}R]$ 在 $T$ 上的一个对象，而
$g$ 是 $x$ 在 $T$ 上其纤维范畴中的一个自同构。换言之，
$x : T \to U$ 且 $g : T \to R$，满足
$x = s \circ g = t \circ g$。这恰好意味着 $g : T \to G$。
惯性纤维化范畴中从 $(x, g) \to (y, h)$ 的一个态射（在 $T$ 上）由
$r : T \to R$ 给出，满足 $s(r) = x$、$t(r) = y$ 且
$c(r, g) = c(h, r)$；参见范畴，引理
\ref{categories-lemma-inertia-fibred-category} 中的交换图表。写成公式为
$$
h = c(r, c(g, i(r))) = c(c(r, g), i(r)).
$$
% P09-ERR-0930
记号 $s(r)$ 等是 $s \circ r$ 等的简写。
% P09-ERR-0929
$r_1 : (x_2, g_2) \to (x_1, g_1)$ 与
$r_2 : (x_1, g_1) \to (x_2, g_2)$ 的复合是
$c(r_1, r_2) : (x_1, g_1) \to (x_3, g_3)$。

\medskip\noindent
注意，在上面，可以用 $g$ 而不用 $(x, g)$ 表示
$\mathcal{I}_{[U/_{\!p}R]}$ 在 $T$ 上的一个对象，因为 $x$ 是 $g$
在结构态射 $G \to U$ 下的像。于是，态射 $g \to h$ 若视为
$\mathcal{I}_{[U/_{\!p}R]}$ 在 $T$ 上的态射，则恰好对应于态射
$r' : T \to R'$，满足
$s'(r') = g$ 且 $t'(r') = h$。此外，复合对应于 (3) 中所述的法则。
引理得证。
\end{proof}

\begin{lemma}
\label{spaces-groupoids-lemma-2-cartesian-inertia}
假设 $B \to S$ 以及 $(U, R, s, t, c)$ 如
定义 \ref{spaces-groupoids-definition-quotient-stack} (1) 中所述。设 $G/U$ 为群胚
$(U, R, s, t, c, e, i)$ 的稳定子群代数空间；参见
定义 \ref{spaces-groupoids-definition-stabilizer-groupoid}。有一个典范 $2$-笛卡儿图表
$$
\xymatrix{
\mathcal{S}_G \ar[r] \ar[d] & \mathcal{S}_U \ar[d] \\
\mathcal{I}_{[U/R]} \ar[r] & [U/R]
}
$$
它由 $(\Sch/S)_{fppf}$ 上的群胚叠组成。
\end{lemma}

\begin{proof}
由引理 \ref{spaces-groupoids-lemma-criterion-fibre-product}，只需证明
% P09-ERR-0931
引理 \ref{spaces-groupoids-lemma-presentation-inertia} 的态射 $s' : R' \to G$
同构于 $s$ 沿结构态射 $G \to U$ 的基变换。这个基变换性质由
$s'$ 的构造显然可见。
\end{proof}

\section{Gerbe 与商叠}
\label{spaces-groupoids-section-gerbes}

\noindent
本节把商叠与 Stacks，第 \ref{stacks-section-gerbes} 节中的讨论联系起来，
特别是与 Stacks，定义
\ref{stacks-definition-gerbe-over-stack-in-groupoids} 中定义的 gerbe 联系起来。
一般而言，本节出现的群胚叠并不是代数叠！

\begin{lemma}
\label{spaces-groupoids-lemma-when-gerbe}
记号与假设同引理 \ref{spaces-groupoids-lemma-quotient-stack-functorial}。商叠的态射
$$
[f] : [U/R] \longrightarrow [U'/R']
$$
$[U/R]$ 成为 $[U'/R']$ 上 gerbe 的充分条件是 $f : U \to U'$ 与
$R \to R'|_U$ 均为 fppf 层的满射。这里 $R'|_U$ 表示 $R'$ 到 $U$
的限制，它是经 $f : U \to U'$ 得到的。
\end{lemma}

\begin{proof}
我们将验证 Stacks，引理 \ref{stacks-lemma-when-gerbe} 的性质
(2) (a) 与 (2) (b) 成立。性质 (2)(a) 成立，因为 $U \to U'$ 是层的满射
（利用引理 \ref{spaces-groupoids-lemma-quotient-stack-objects} 可知，$[U'/R']$ 中的对象局部来自
$U'$）。为证明 (2)(b)，令 $x, y$ 为 $[U/R]$ 的对象，它们位于概形 $T/S$ 上。
% P09-ERR-0932
令 $x', y'$ 为 $x, y$ 在范畴 $[U'/'R]_T$ 中的像。条件 (2)(b) 要求我们检验层的映射
$$
\mathit{Isom}(x, y) \longrightarrow \mathit{Isom}(x', y')
$$
在 $(\Sch/T)_{fppf}$ 上是满射。为此，可以在 $T$ 上 fppf 局部地工作，
% P09-ERR-0933
并假设它们来自 $a, b \in U(T)$。在这种情形下，$x', y'$ 分别对应于
$f \circ a, f \circ b$。由引理 \ref{spaces-groupoids-lemma-quotient-stack-morphisms}，
在这种情形下，上述层的映射变为
$$
T \times_{(a, b), U \times_B U} R
\longrightarrow
T \times_{f \circ a, f \circ b, U' \times_B U'} R' =
T \times_{(a, b), U \times_B U} R'|_U.
$$
因此，$R \to R'|_U$ 是 $(\Sch/S)_{fppf}$ 上 fppf 层的满射这一假设
蕴含所需的满射性。
\end{proof}

\begin{lemma}
\label{spaces-groupoids-lemma-group-quotient-gerbe}
设 $S$ 为概形，$B$ 为 $S$ 上的代数空间，$G$ 为 $B$ 上的群代数空间。
赋予 $B$ 以 $G$ 的平凡作用。态射
$$
[B/G] \longrightarrow \mathcal{S}_B
$$
（引理 \ref{spaces-groupoids-lemma-quotient-stack-arrows}）使 $[B/G]$ 成为 $B$ 上的 gerbe。
\end{lemma}

\begin{proof}
由引理 \ref{spaces-groupoids-lemma-when-gerbe} 立即可得，因为态射 $B \to B$ 与
$B \times_B G \to B$ 作为层的态射都是满射。
\end{proof}

\section{商叠与大位点的变换}
\label{spaces-groupoids-section-bigger-site}

\noindent
我们建议初次阅读时跳过本节。叠的拉回定义于
Stacks，第 \ref{stacks-section-inverse-image} 节。

\begin{lemma}
\label{spaces-groupoids-lemma-quotient-stack-change-big-site}
设给定大位点 $\Sch_{fppf}$ 与 $\Sch'_{fppf}$。假设
$\Sch_{fppf}$ 包含于 $\Sch'_{fppf}$ 中；参见拓扑，第
\ref{topologies-section-change-alpha} 节。设
$S \in \Ob(\Sch_{fppf})$。设
$B, U, R \in \Sh((\Sch/S)_{fppf})$ 为代数空间，并设
$(U, R, s, t, c)$ 为 $B$ 上代数空间中的群胚。
% P09-ERR-0934
设 $f : (\Sch'/S)_{fppf} \to (\Sch/S)_{fppf}$ 为与包含函子
$u : \Sch_{fppf} \to \Sch'_{fppf}$ 对应的位点态射。则有典范等价
$$
[f^{-1}U/f^{-1}R]
\longrightarrow
f^{-1}[U/R]
$$
它由 $(\Sch'/S)_{fppf}$ 上的群胚叠组成。
\end{lemma}

\begin{proof}
注意，$f^{-1}B, f^{-1}U, f^{-1}R \in \Sh((\Sch'/S)_{fppf})$
由空间，引理 \ref{spaces-lemma-change-big-site} 是代数空间，因而
$(f^{-1}U, f^{-1}R, f^{-1}s, f^{-1}t, f^{-1}c)$ 是
$f^{-1}B$ 上代数空间中的群胚。因此该陈述有意义。

\medskip\noindent
% P09-ERR-0935
范畴 $u_p[U/_{\!p}R]$ 是范畴 $u_{pp}[U/_{\!p}R]$ 关于态射的
右乘法系 $I$ 的局部化。$u_{pp}[U/_{\!p}R]$ 的一个对象是三元组
$$
(T', \phi : T' \to T, x)
$$
其中
$T' \in \Ob((\Sch'/S)_{fppf})$，
$T \in \Ob((\Sch/S)_{fppf})$，$\phi$ 是 $S$ 上概形的态射，
而 $x : T \to U$ 是 $(\Sch/S)_{fppf}$ 上的层态射。注意，概形态射
$\phi : T' \to T$ 等同于态射 $\phi : T' \to u(T)$；又由于
$u(T)$ 表示 $f^{-1}T$，这也等同于态射 $T' \to f^{-1}T$。
此外，由于 $f^{-1}$ 在代数空间上全忠实（参见空间，引理
\ref{spaces-lemma-fully-faithful}），我们也可以把 $x$ 视为态射
$x : f^{-1}T \to f^{-1}U$。下文将不再说明地作这些等同。态射
$$
(a, a', \alpha) :
(T'_1, \phi_1 : T'_1 \to T_1, x_1)
\longrightarrow
(T'_2, \phi_2 : T'_2 \to T_2, x_2)
$$
若是 $u_{pp}[U/_{\!p}R]$ 的态射，就是一个交换图表
$$
\xymatrix{
& & U \\
T'_1 \ar[d]_{a'} \ar[r]_{\phi_1} &
T_1 \ar[d]_a \ar[ru]^{x_1} \ar[r]_\alpha &
R \ar[d]^t \ar[u]_s \\
T'_2 \ar[r]^{\phi_2} &
T_2 \ar[r]^{x_2} &
U
}
$$
而且这样的态射属于 $I$ 当且仅当 $T'_1 = T'_2$ 且
$a' = \text{id}$。我们按如下法则定义函子
$$
u_{pp}[U/_{\!p}R] \longrightarrow [f^{-1}U/_{\!p}f^{-1}R]
$$
它在对象上的法则为
$$
(T', \phi : T' \to T, x) \longmapsto (x \circ \phi : T' \to f^{-1}U)
$$
在上述态射上的法则为
$$
(a, a', \alpha) \longmapsto (\alpha \circ \phi_1 : T'_1 \to f^{-1}R)
$$
显然，$I$ 的元素被变换为同构，因为
$(f^{-1}U, f^{-1}R, f^{-1}s, f^{-1}t, f^{-1}c)$ 是
$f^{-1}B$ 上代数空间中的群胚。因此，这个函子典范地通过函子
$$
u_p[U/_{\!p}R] \longrightarrow [f^{-1}U/_{\!p}f^{-1}R]
$$
分解。应用叠化，可得叠函子
$$
f^{-1}[U/R] \longrightarrow [f^{-1}U/f^{-1}R]
$$
它在 $(\Sch'/S)_{fppf}$ 上；这是因为由 Stacks，引理
\ref{stacks-lemma-technical-up}，
叠 $f^{-1}[U/R]$ 是 $u_p[U/_{\!p}R]$ 的叠化。

\medskip\noindent
至此我们已有叠的态射。要验证它是等价，只需证明它全忠实且对象局部地
位于本质像中；参见 Stacks，引理 \ref{stacks-lemma-characterize-ff} 与
\ref{stacks-lemma-characterize-essentially-surjective-when-ff}。
% P09-ERR-0936
关于对象的陈述成立，因为 $f^{-1}R$ 有满平展态射
$f^{-1}W \to f^{-1}R$，其中 $W$ 是 $(\Sch/S)_{fppf}$ 的某个对象。
要证明函子“满”，只需证明态射局部地位于该函子的像中；这成立是因为
$f^{-1}U$ 有满平展态射 $f^{-1}W \to f^{-1}U$，其中 $W$ 是
$(\Sch/S)_{fppf}$ 的某个对象。略去函子忠实性的证明。
\end{proof}

\section{分离性条件}
\label{spaces-groupoids-section-separation}

\noindent
这实际上意指态射 $j : R \to U \times_B U$ 所满足的条件；这里给定一个
代数空间中的群胚 $(U, R, s, t, c)$，它在 $B$ 上。与上一节一样，
我们先写出相应的图表。

\begin{lemma}
\label{spaces-groupoids-lemma-diagram-diagonal}
设 $B \to S$ 如第 \ref{spaces-groupoids-section-notation} 节所示。设
$(U, R, s, t, c)$ 为 $B$ 上代数空间中的群胚。设
$G \to U$ 为稳定子群代数空间。交换图表
$$
\xymatrix{
R \ar[d]^{\Delta_{R/U \times_B U}} \ar[rrr]_{f \mapsto (f, s(f))} & & &
R \times_{s, U} U \ar[d] \ar[r] & U \ar[d] \\
R \times_{(U \times_B U)} R \ar[rrr]^{(f, g) \mapsto (f, f^{-1} \circ g)} & & &
R \times_{s, U} G \ar[r] & G
}
$$
% P09-ERR-0937
中的左边两条水平箭头是同构，而右方块是纤维积方块。
\end{lemma}

\begin{proof}
略。这是关于定义以及代数几何中函子观点的练习。
\end{proof}

\begin{lemma}
\label{spaces-groupoids-lemma-diagonal}
设 $B \to S$ 如第 \ref{spaces-groupoids-section-notation} 节所示。设
$(U, R, s, t, c)$ 为 $B$ 上代数空间中的群胚。设
$G \to U$ 为稳定子群代数空间。
\begin{enumerate}
\item 下列条件等价：
\begin{enumerate}
\item $j : R \to U \times_B U$ 是分离的；
\item $G \to U$ 是分离的；以及
\item $e : U \to G$ 是闭浸入。
\end{enumerate}
\item 下列条件等价：
\begin{enumerate}
\item $j : R \to U \times_B U$ 是局部分离的；
\item $G \to U$ 是局部分离的；以及
\item $e : U \to G$ 是浸入。
\end{enumerate}
\item 下列条件等价：
\begin{enumerate}
\item $j : R \to U \times_B U$ 是拟分离的；
\item $G \to U$ 是拟分离的；以及
\item $e : U \to G$ 是拟紧的。
\end{enumerate}
\end{enumerate}
\end{lemma}

\begin{proof}
群代数空间 $G \to U$ 是 $R \to U \times_B U$ 沿对角态射
$U \to U \times_B U$ 的基变换；参见引理
\ref{spaces-groupoids-lemma-groupoid-stabilizer}。因此，如果 $j$ 是分离的（相应地，
局部分离的、拟分离的），那么 $G \to U$ 是分离的（相应地，
局部分离的、拟分离的）。参见空间的态射，引理
\ref{spaces-morphisms-lemma-base-change-separated}。因此在 (1)、(2) 与
(3) 中，(a) $\Rightarrow$ (b)。

\medskip\noindent
反之，如果 $G \to U$ 是分离的（相应地，局部分离的、拟分离的），
那么态射 $e : U \to G$ 作为结构态射 $G \to U$ 的截面，是闭浸入
（相应地，浸入、拟紧的）；参见空间的态射，引理
\ref{spaces-morphisms-lemma-section-immersion}。因此在 (1)、(2) 与
(3) 中，(b) $\Rightarrow$ (c)。

\medskip\noindent
如果 $e$ 是闭浸入（相应地，浸入、拟紧的），那么由引理
\ref{spaces-groupoids-lemma-diagram-diagonal} 的结论（以及空间，引理
\ref{spaces-lemma-base-change-immersions} 和空间的态射，引理
\ref{spaces-morphisms-lemma-base-change-quasi-compact}），可知
$\Delta_{R/U \times_B U}$ 是闭浸入（相应地，浸入、拟紧的）。
因此在 (1)、(2) 与 (3) 中，(c) $\Rightarrow$ (a)。
\end{proof}
